% \documentclass[lettersize,journal]{IEEEtran}
\documentclass[10pt,journal,compsoc,x11names]{IEEEtran}
\usepackage{amsmath,amsfonts}
\usepackage{algorithmic}
\usepackage{algorithm}
\usepackage{array}
\usepackage[caption=false,font=normalsize,labelfont=sf,textfont=sf]{subfig}
\usepackage{textcomp}
\usepackage{stfloats}
\usepackage{url}
\usepackage{verbatim}
\usepackage{graphicx}
\usepackage{cite}
\usepackage{multirow}
\hyphenation{op-tical net-works semi-conduc-tor IEEE-Xplore}
% updated with editorial comments 8/9/2021

\usepackage{hyperref}
\hypersetup{
    colorlinks=true,      % 启用彩色链接
    linkcolor=blue,       % 设置章节、图表等的链接颜色
    citecolor=red,        % 设置参考文献的链接颜色
    filecolor=magenta,    % 设置文件链接的颜色
    urlcolor=blue        % 设置网址链接的颜色
}


\usepackage{forest}
\usetikzlibrary{shadows}
\usepackage{amsmath}
\let\Bbbk\relax
\usepackage{amssymb}
\usepackage{cleveref}

\definecolor{deepred}{rgb}{0.698,0.133,0.133}
\definecolor{blue}{rgb}{0,0,1}
\usepackage{tikz}
\usepackage{amssymb}
\newcommand\encircle[2][]{\tikz[overlay]\node[fill=blue!20,inner sep=2pt, anchor=text, rectangle, rounded corners=1.5mm,#1] {#2};\phantom{#2}}

\definecolor{lightcoral}{rgb}{0.94, 0.5, 0.5}
\definecolor{lightgreen}{rgb}{0.56, 0.93, 0.56}
\definecolor{harvestgold}{rgb}{0.85, 0.57, 0.0}
\definecolor{brightlavender}{rgb}{0.75, 0.58, 0.89}
\definecolor{capri}{rgb}{0.0, 0.75, 1.0}
\definecolor{carminepink}{rgb}{0.92, 0.3, 0.26}
\definecolor{celadon}{rgb}{0.67, 0.88, 0.69}
\definecolor{darkpastelgreen}{rgb}{0.01, 0.75, 0.24}

\definecolor{deepred}{rgb}{0.698,0.133,0.133}
\definecolor{blue}{rgb}{0,0,1}
\definecolor{lightergray}{RGB}{230,230,230}
\definecolor{DarkRed}{RGB}{130,25,0}
\definecolor{PurpleRed}{RGB}{204,0,102}
\definecolor{DarkGreen}{RGB}{30,130,30}
\definecolor{DarkBlue}{RGB}{0,0,250}
\definecolor{DarkYellow}{RGB}{255,128,0}
\definecolor{light-gray}{gray}{0.95}
\definecolor{lightgreen}{RGB}{231,255,219}
\definecolor{lightred}{RGB}{252,231,234}
\definecolor{lightyellow}{RGB}{250,253,191}
\definecolor{lightpurple}{RGB}{229,204,255}
\definecolor{lightblue}{RGB}{229,246,254}
\definecolor{value-modification}{RGB}{250, 217, 86}
\definecolor{digit-expansion}{RGB}{216, 194, 104}
\definecolor{integer-decimal-fraction}{RGB}{240, 133, 51}
\definecolor{semantic-paraphrasing}{RGB}{85, 157, 63}
\definecolor{complexity-increasing}{RGB}{58, 120, 175}
\definecolor{question-transformation}{RGB}{174, 205, 225}
\definecolor{interference-injection}{RGB}{255,204,229}
\definecolor{remove-constrain}{RGB}{204,204,255}
\definecolor{myGreen}{RGB}{127,210,85}
\definecolor{myOrange}{RGB}{242,154,66}
\definecolor{myYellow}{RGB}{247,223,65}
\definecolor{myRed}{RGB}{232,80,43}
\definecolor{myViolet}{RGB}{162,57,102}
\definecolor{myBlue}{HTML}{4686f3}
\definecolor{myYellowv2}{HTML}{E6C802}
\definecolor{myOrangev2}{HTML}{ED8E55}
\definecolor{MyGreenv2}{HTML}{009B55}
\definecolor{MyRedv2}{HTML}{c22f2f}

\newcommand{\redcross}{{\color{red}\ding{55}}} % 定义红叉命令
% new commands
\definecolor{lightcoral}{rgb}{0.94, 0.5, 0.5}
\definecolor{lightgreen}{rgb}{0.56, 0.93, 0.56}
\definecolor{harvestgold}{rgb}{0.85, 0.57, 0.0}
\definecolor{brightlavender}{rgb}{0.75, 0.58, 0.89}
\definecolor{capri}{rgb}{0.0, 0.75, 1.0}
\definecolor{carminepink}{rgb}{0.92, 0.3, 0.26}
\definecolor{celadon}{rgb}{0.67, 0.88, 0.69}
\definecolor{darkpastelgreen}{rgb}{0.01, 0.75, 0.24}

\definecolor{deepred}{rgb}{0.698,0.133,0.133}
\definecolor{blue}{rgb}{0,0,1}

% \newcommand{\redcross}{{\color{red}\ding{55}}} % 定义红叉命令

\newcommand*\circled[1]{\tikz[baseline=(char.base)]{
            \node[shape=circle,draw,inner sep=0.5pt,thick] (char) {#1};}}


\crefformat{section}{\S#2#1#3} 
\crefformat{subsection}{\S#2#1#3}
\crefformat{subsubsection}{\S#2#1#3}

% Math
\usepackage{amsthm}
\usepackage{amsmath}%
\usepackage{MnSymbol}%
% \usepackage{wasysym}%
\newcommand\smallO{
  \mathchoice
    {{\scriptstyle\mathcal{O}}}% \displaystyle
    {{\scriptstyle\mathcal{O}}}% \textstyle
    {{\scriptscriptstyle\mathcal{O}}}% \scriptstyle
    {\scalebox{.7}{$\scriptscriptstyle\mathcal{O}$}}%\scriptscriptstyle
  }
\makeatletter
\newtheorem*{rep@theorem}{\rep@title}
\newcommand{\newreptheorem}[2]{%
\newenvironment{rep#1}[1]{%
 \def\rep@title{#2 \ref{##1}}%
 \begin{rep@theorem}}%
 {\end{rep@theorem}}}
\makeatother
\newtheorem{theorem}{Theorem}[section]
\newreptheorem{theorem}{Theorem}
\newtheorem{corollary}[theorem]{Corollary}
\newtheorem{lemma}[theorem]{Lemma}
\newreptheorem{lemma}{Lemma}
\newtheorem{proposition}[theorem]{Proposition}
\newreptheorem{proposition}{proposition}
\newtheorem{definition}[theorem]{Definition}
\newtheorem{assumption}[theorem]{Assumption}
\newtheorem{remark}[theorem]{Remark}
\usepackage{bm}

% Tables
\usepackage{booktabs}
\usepackage{xcolor,pifont}
\usepackage{bbding}
\usepackage{colortbl}

% List
\usepackage[inline,shortlabels]{enumitem}

% ORCID
\usepackage{tikz,xcolor,hyperref}% Make Orcid icon
\definecolor{lime}{HTML}{A6CE39}
\DeclareRobustCommand{\orcidicon}{%
    \begin{tikzpicture}
    \draw[lime, fill=lime] (0,0) 
    circle [radius=0.16] 
    node[white] {{\fontfamily{qag}\selectfont \tiny ID}};    \draw[white, fill=white] (-0.0625,0.095) 
    circle [radius=0.007];    \end{tikzpicture}
    \hspace{-2mm}}
\foreach \x in {A, ..., Z}{%
    \expandafter\xdef\csname orcid\x\endcsname{\noexpand\href{https://orcid.org/\csname orcidauthor\x\endcsname}{\noexpand\orcidicon}}
}

% Define the ORCID iD command for each author separately.
\newcommand{\orcidauthorA}{0000-0002-6743-4175}



\begin{document}
% \title{Fast to Slow: From {System 1} to {System 2} Reasoning in Large Language Models}

\title{From System 1 to System 2: A Survey of\\ Reasoning Large Language Models}

\author{
Zhong-Zhi~Li$^{*}$,
Duzhen~Zhang$^{*}$,
Ming-Liang~Zhang$^{\S}$,
Jiaxin~Zhang$^{\S}$,
Zengyan~Liu$^{\S}$,
Yuxuan~Yao$^{\S}$,\\
Haotian~Xu, 
Junhao~Zheng, 
Pei-Jie Wang, 
Xiuyi~Chen,
Yingying~Zhang,
Fei~Yin, 
Jiahua~Dong, \\
Zhijiang~Guo$^{\dagger}$,
Le~Song$^{\dagger}$,
Cheng-Lin~Liu$^{\dagger}$\orcidA{},~\IEEEmembership{Fellow,~IEEE}
\thanks{Version: v1 (major update on February 24, 2025)}
\thanks{$^{*}$Core contribution. $^{\S}$Significant contribution. $^\dagger$Corresponding author.}
\thanks{Duzhen Zhang, Jiahua Dong, and Le Song are with the Mohamed bin Zayed University of Artificial Intelligence, Abu Dhabi, UAE (E-mail: bladedancer957@gmail.com; dongjiahua1995@gmail.com; le.song@mbzuai.ac.ae).}
\thanks{Zhong-Zhi Li, Pei-Jie Wang, Xiuyi Chen, Fei Yin, and Cheng-Lin Liu are with the Institute of Automation, Chinese Academy of Sciences, Beijing, China (E-mail: lizhongzhi2022@ia.ac.cn; hugheren.chan@gmail.com; fyin@nlpr.ia.ac.cn; liucl@nlpr.ia.ac.cn).}
\thanks{Ming-Liang Zhang is with the AiShiWeiLai AI Research, Beijing, China (E-mail: zhangmingliang@yuaiweiwu.com).}
% \thanks{Zenyan Liu and Yuxuan Yao are with the City University of Hong Kong, Hong Kong, China (E-mail: zengyaliu2-c@my.cityu.edu.hk; yuxuanyao3-c@my.cityu.edu.hk).}
\thanks{Zengyan Liu, Yuxuan Yao, and Zhijiang Guo is with City University of Hong Kong and the Hong Kong University of Science and Technology (Guangzhou), China (E-mail: zengyaliu2-c@my.cityu.edu.hk; yuxuanyao3-c@my.cityu.edu.hk; zhijiangguo@hkust-gz.edu.cn).}
\thanks{Jiaxin Zhang is with the University of Strathclyde, Glasgow, UK (E-mail: jiaxin.zhang@strath.ac.uk).}
\thanks{Haotian Xu is with the Xiaohongshu Inc, Beijing, China (E-mail: xuhaotian@xiaohongshu.com).}
\thanks{Yingying Zhang is with the East China Normal University, Shanghai, China (E-mail: yyzhang@fem.ecnu.edu.cn).}
\thanks{Junhao Zheng is with the South China University of Technology, Guangzhou, China (E-mail: junhaozheng47@outlook.com).}
}% <-this % stops a space
% \thanks{Manuscript received April 19, 2021; revised August 16, 2021.}}

% The paper headers
\markboth{Journal of \LaTeX\ Class Files, January 2025}%
{Shell \MakeLowercase{\textit{et al.}}: A Sample Article Using IEEEtran.cls for IEEE Journals}

% \IEEEpubid{0000--0000/00\$00.00~\copyright~2021 IEEE}
% Remember, if you use this you must call \IEEEpubidadjcol in the second
% column for its text to clear the IEEEpubid mark.


\IEEEtitleabstractindextext{

\begin{abstract}  
Test time scaling is currently one of the most active research areas that shows promise after training time scaling has reached its limits.
Deep-thinking (DT) models are a class of recurrent models that can perform easy-to-hard generalization by assigning more compute to harder test samples.
However, due to their inability to determine the complexity of a test sample, DT models have to use a large amount of computation for both easy and hard test samples.
Excessive test time computation is wasteful and can cause the ``overthinking'' problem where more test time computation leads to worse results.
In this paper, we introduce a test time training method for determining the optimal amount of computation needed for each sample during test time.
We also propose Conv-LiGRU, a novel recurrent architecture for efficient and robust visual reasoning. 
Extensive experiments demonstrate that Conv-LiGRU is more stable than DT, effectively mitigates the ``overthinking'' phenomenon, and achieves superior accuracy.
\end{abstract}  

\begin{IEEEkeywords}
Slow-thinking, Large Language Models, Human-like Reasoning, Decision Making in AI, AGI
\end{IEEEkeywords}
}

\maketitle



% 1. 主线 system1 to system2  可以讲 近十几年 从dnn 到llm 到o1； 



% 2. 然后到o1的这个过程 其实是 融合了之前的一些slow-thinking思想，，，介绍一下这些slow思想的发展过程，以及和现在o1技术的联系，



% 3. slow发展其实是 纯slow的符号逻辑推理，后来是概率图， rl，后来dnn兴起后，slow其实和dnn有过结合，比如神经符号，gnn，深度rl，


% 4. o1 算是 slow思想和llm结合的高级slow， 和之前slow+dnn不同的是，，这个更open，之前主要是close 或者specific tasks




\section{Introduction}
\label{sec:introduction}
The business processes of organizations are experiencing ever-increasing complexity due to the large amount of data, high number of users, and high-tech devices involved \cite{martin2021pmopportunitieschallenges, beerepoot2023biggestbpmproblems}. This complexity may cause business processes to deviate from normal control flow due to unforeseen and disruptive anomalies \cite{adams2023proceddsriftdetection}. These control-flow anomalies manifest as unknown, skipped, and wrongly-ordered activities in the traces of event logs monitored from the execution of business processes \cite{ko2023adsystematicreview}. For the sake of clarity, let us consider an illustrative example of such anomalies. Figure \ref{FP_ANOMALIES} shows a so-called event log footprint, which captures the control flow relations of four activities of a hypothetical event log. In particular, this footprint captures the control-flow relations between activities \texttt{a}, \texttt{b}, \texttt{c} and \texttt{d}. These are the causal ($\rightarrow$) relation, concurrent ($\parallel$) relation, and other ($\#$) relations such as exclusivity or non-local dependency \cite{aalst2022pmhandbook}. In addition, on the right are six traces, of which five exhibit skipped, wrongly-ordered and unknown control-flow anomalies. For example, $\langle$\texttt{a b d}$\rangle$ has a skipped activity, which is \texttt{c}. Because of this skipped activity, the control-flow relation \texttt{b}$\,\#\,$\texttt{d} is violated, since \texttt{d} directly follows \texttt{b} in the anomalous trace.
\begin{figure}[!t]
\centering
\includegraphics[width=0.9\columnwidth]{images/FP_ANOMALIES.png}
\caption{An example event log footprint with six traces, of which five exhibit control-flow anomalies.}
\label{FP_ANOMALIES}
\end{figure}

\subsection{Control-flow anomaly detection}
Control-flow anomaly detection techniques aim to characterize the normal control flow from event logs and verify whether these deviations occur in new event logs \cite{ko2023adsystematicreview}. To develop control-flow anomaly detection techniques, \revision{process mining} has seen widespread adoption owing to process discovery and \revision{conformance checking}. On the one hand, process discovery is a set of algorithms that encode control-flow relations as a set of model elements and constraints according to a given modeling formalism \cite{aalst2022pmhandbook}; hereafter, we refer to the Petri net, a widespread modeling formalism. On the other hand, \revision{conformance checking} is an explainable set of algorithms that allows linking any deviations with the reference Petri net and providing the fitness measure, namely a measure of how much the Petri net fits the new event log \cite{aalst2022pmhandbook}. Many control-flow anomaly detection techniques based on \revision{conformance checking} (hereafter, \revision{conformance checking}-based techniques) use the fitness measure to determine whether an event log is anomalous \cite{bezerra2009pmad, bezerra2013adlogspais, myers2018icsadpm, pecchia2020applicationfailuresanalysispm}. 

The scientific literature also includes many \revision{conformance checking}-independent techniques for control-flow anomaly detection that combine specific types of trace encodings with machine/deep learning \cite{ko2023adsystematicreview, tavares2023pmtraceencoding}. Whereas these techniques are very effective, their explainability is challenging due to both the type of trace encoding employed and the machine/deep learning model used \cite{rawal2022trustworthyaiadvances,li2023explainablead}. Hence, in the following, we focus on the shortcomings of \revision{conformance checking}-based techniques to investigate whether it is possible to support the development of competitive control-flow anomaly detection techniques while maintaining the explainable nature of \revision{conformance checking}.
\begin{figure}[!t]
\centering
\includegraphics[width=\columnwidth]{images/HIGH_LEVEL_VIEW.png}
\caption{A high-level view of the proposed framework for combining \revision{process mining}-based feature extraction with dimensionality reduction for control-flow anomaly detection.}
\label{HIGH_LEVEL_VIEW}
\end{figure}

\subsection{Shortcomings of \revision{conformance checking}-based techniques}
Unfortunately, the detection effectiveness of \revision{conformance checking}-based techniques is affected by noisy data and low-quality Petri nets, which may be due to human errors in the modeling process or representational bias of process discovery algorithms \cite{bezerra2013adlogspais, pecchia2020applicationfailuresanalysispm, aalst2016pm}. Specifically, on the one hand, noisy data may introduce infrequent and deceptive control-flow relations that may result in inconsistent fitness measures, whereas, on the other hand, checking event logs against a low-quality Petri net could lead to an unreliable distribution of fitness measures. Nonetheless, such Petri nets can still be used as references to obtain insightful information for \revision{process mining}-based feature extraction, supporting the development of competitive and explainable \revision{conformance checking}-based techniques for control-flow anomaly detection despite the problems above. For example, a few works outline that token-based \revision{conformance checking} can be used for \revision{process mining}-based feature extraction to build tabular data and develop effective \revision{conformance checking}-based techniques for control-flow anomaly detection \cite{singh2022lapmsh, debenedictis2023dtadiiot}. However, to the best of our knowledge, the scientific literature lacks a structured proposal for \revision{process mining}-based feature extraction using the state-of-the-art \revision{conformance checking} variant, namely alignment-based \revision{conformance checking}.

\subsection{Contributions}
We propose a novel \revision{process mining}-based feature extraction approach with alignment-based \revision{conformance checking}. This variant aligns the deviating control flow with a reference Petri net; the resulting alignment can be inspected to extract additional statistics such as the number of times a given activity caused mismatches \cite{aalst2022pmhandbook}. We integrate this approach into a flexible and explainable framework for developing techniques for control-flow anomaly detection. The framework combines \revision{process mining}-based feature extraction and dimensionality reduction to handle high-dimensional feature sets, achieve detection effectiveness, and support explainability. Notably, in addition to our proposed \revision{process mining}-based feature extraction approach, the framework allows employing other approaches, enabling a fair comparison of multiple \revision{conformance checking}-based and \revision{conformance checking}-independent techniques for control-flow anomaly detection. Figure \ref{HIGH_LEVEL_VIEW} shows a high-level view of the framework. Business processes are monitored, and event logs obtained from the database of information systems. Subsequently, \revision{process mining}-based feature extraction is applied to these event logs and tabular data input to dimensionality reduction to identify control-flow anomalies. We apply several \revision{conformance checking}-based and \revision{conformance checking}-independent framework techniques to publicly available datasets, simulated data of a case study from railways, and real-world data of a case study from healthcare. We show that the framework techniques implementing our approach outperform the baseline \revision{conformance checking}-based techniques while maintaining the explainable nature of \revision{conformance checking}.

In summary, the contributions of this paper are as follows.
\begin{itemize}
    \item{
        A novel \revision{process mining}-based feature extraction approach to support the development of competitive and explainable \revision{conformance checking}-based techniques for control-flow anomaly detection.
    }
    \item{
        A flexible and explainable framework for developing techniques for control-flow anomaly detection using \revision{process mining}-based feature extraction and dimensionality reduction.
    }
    \item{
        Application to synthetic and real-world datasets of several \revision{conformance checking}-based and \revision{conformance checking}-independent framework techniques, evaluating their detection effectiveness and explainability.
    }
\end{itemize}

The rest of the paper is organized as follows.
\begin{itemize}
    \item Section \ref{sec:related_work} reviews the existing techniques for control-flow anomaly detection, categorizing them into \revision{conformance checking}-based and \revision{conformance checking}-independent techniques.
    \item Section \ref{sec:abccfe} provides the preliminaries of \revision{process mining} to establish the notation used throughout the paper, and delves into the details of the proposed \revision{process mining}-based feature extraction approach with alignment-based \revision{conformance checking}.
    \item Section \ref{sec:framework} describes the framework for developing \revision{conformance checking}-based and \revision{conformance checking}-independent techniques for control-flow anomaly detection that combine \revision{process mining}-based feature extraction and dimensionality reduction.
    \item Section \ref{sec:evaluation} presents the experiments conducted with multiple framework and baseline techniques using data from publicly available datasets and case studies.
    \item Section \ref{sec:conclusions} draws the conclusions and presents future work.
\end{itemize}

\section{Foundations of Reasoning LLMs}\label{early_srs}

% rl supervise, prm->rl /reward shaping技术, macro-action->symbolic模板或者专家规则,  struture mcts-> search/概率图，self-improve alphago


In this section, we provide a concise overview of the progress in foundational LLMs and the early development of key \textit{System 2} technologies, highlighting critical advancements that, when combined with foundational LLMs, have paved the way for reasoning LLMs. These advancements include symbolic logic systems, MCTS, and RL.


\subsection{Foundational LLMs}\label{f_llm}


The development of foundational LLMs saw significant advancements with the introduction of pretrained Transformers \cite{vaswani2017attention} in 2018-2019, notably through BERT \cite{DBLP:conf/naacl/DevlinCLT19} and GPT \cite{radford2018improving}. 
These models leveraged unsupervised pretraining on vast text corpora, followed by fine-tuning for task-specific applications. This approach enabled them to develop a broad language understanding before specializing in tasks such as sentiment analysis, entity recognition, and question answering. 
BERT's bidirectional context processing improved word understanding, while GPT excelled in text generation with its unidirectional design.

The release of GPT-2 \cite{radford2019language} in 2019, with 1.5 billion parameters, marked a significant leap in generative performance, though it also raised ethical concerns. 
GPT-3 \cite{brown2020language}, with 175 billion parameters, further demonstrated the power of unsupervised pretraining, excelling in few-shot learning and performing well across a wide range of NLP tasks. 
In subsequent years, multimodal models like CLIP \cite{radford2021learning} and DALL-E \cite{ramesh2021zero} emerged, integrating text and visual inputs. 
These models enabled new tasks, such as generating images from text, and enhanced human-computer interaction.


By 2023-2024, models such as GPT-4/4o \cite{openai2023gpt4,gpt4o-0513}, LLaMA \cite{touvron2023llama}, and LLaVA \cite{liu2023llava} demonstrated advanced capabilities in reasoning, contextual understanding, and multimodal reasoning, processing both text and images \cite{alayrac2022flamingo, DBLP:conf/icml/0008LSH23, DBLP:journals/corr/abs-2305-06500}. 
DeepSeek-V3 \cite{liu2024deepseek}, featuring a 671B Mixture-of-Expert architecture \cite{he2021fastmoe, du2022glam, dai2024deepseekmoe}, outperforms several other LLMs on key benchmarks while offering significant improvements in efficiency and processing speed. 
The evolution of foundational LLMs has revolutionized AI, enabling more sophisticated applications in language comprehension, problem-solving, and human-machine collaboration.



\noindent\textbf{Summary:} The development of foundational LLMs has progressed from pretrained transformers like BERT to multimodal models such as GPT-4, enhancing language understanding, text generation, and image processing. 
This advancement has led to significant breakthroughs in AI, improving language comprehension, problem-solving, and human-computer interaction. Building on deep learning advancements \cite{rumelhart1986learning, lecun1995convolutional, hochreiter1997long, hinton2006fast, hinton2006reducing, hinton2012deep, krizhevsky2012imagenet, DBLP:conf/emnlp/ChoMGBBSB14, sutskever2014sequence, srivastava2014dropout, kingma2014adam, lecun2015deep, he2016deep, huang2017densely,vaswani2017attention, goodfellow2020generative}, foundational LLMs can learn extensive world knowledge and semantic relationships from vast textual or multimodal data. This enables them to exhibit emergent capabilities such as In-Context Learning (ICL) \cite{min2022rethinking, dong2024survey}, prompt engineering \cite{white2023prompt, lester2021power}, and Chain-of-Thought (CoT) reasoning \cite{wei2022chain}, significantly enhancing their adaptability and creative problem-solving abilities.

Despite this progress, foundational LLMs operate similarly to \textit{System 1} reasoning, relying on fast, heuristic-driven decision-making and lacking the step-by-step analysis characteristic of \textit{System 2}. However, their developments lay a solid foundation for future reasoning LLMs\textendash especially when integrated with the following early \textit{System 2} technologies. This combination paves the way for more versatile, flexible, and human-like reasoning models.












\subsection{Symbolic Logic Systems}\label{symb_exp}


Symbolic logic systems mark the earliest phase of AI, utilizing rules and logical principles to represent knowledge and draw conclusions \cite{lewis1959symbolic, carnap2012introduction}. 
They are particularly effective in structured domains, where formal logic ensures precision.

Prolog, a logic programming language based on first-order logic, allows users to define facts, rules, and reason through queries. 
It has been pivotal in symbolic reasoning systems, especially in NLP and expert systems \cite{colmerauer1990introduction, clocksin2003programming, apt1997logic}. Logic-based systems like Prolog employ propositional and predicate logic for formal reasoning \cite{singh1999formal, jeroslow1988computation}. 
From the 1960s to the early 1980s, this approach dominated AI, with systems like IBM's LISP \cite{mccarthy1978history} for symbolic computation and Resolution Theorem Provers \cite{bachmair2001resolution} for automated reasoning. 
In the 1970s, Marvin Minsky introduced Frames, which organized knowledge into structured frameworks, influencing both expert systems and cognitive science \cite{minsky1974framework}.


\noindent\textbf{Summary:} 
Symbolic logic systems were pivotal milestones in early AI development. 
Based on formal logic, they excelled in well-defined problems, particularly in structured environments. 
However, they also exposed the limitations of rigid, rule-based systems. Despite these constraints, symbolic logic remains foundational to the progress of AI.

Recent advancements in reasoning LLMs have greatly enhanced the emulation of human-like \textit{System 2} cognitive processes through sophisticated thought architectures, known as Macro Action frameworks (Section \ref{macro_action}). By combining symbolic templates or rules with foundational LLMs, macro actions have significantly improved their reasoning capabilities. 
Integrating macro actions into foundational LLMs has transformed their ability to handle complex reasoning tasks, as hierarchical planning allows models to make high-level decisions before delving into specific problem details, mirroring symbolic logic's structured approach.

% Furthermore, recent breakthroughs in combining LLMs with symbolic systems have demonstrated their potential. 
% For instance, Google's AlphaGeometry \cite{trinh2024solving} and AlphaGeometry2 \cite{chervonyi2025gold} combine LLMs with symbolic engines, achieving success in the International Olympiad in Mathematics (IMO). 
% These symbolic engines offer clear reasoning pathways, making them effective for tackling complex or ambiguous problems. 
% While early symbolic systems were less adaptable, they provided better explainability and clearer reasoning steps. 
% When combined with LLMs, they lead to more reliable and interpretable performance. 
% This synergy underscores the continued relevance of symbolic approaches in modern AI.




% \subsubsection{Expert Systems}\label{exp}

% XCON is a configuration expert system that helps design custom computer systems by optimizing hardware and software based on customer specifications \cite{}. 
% Its commercial success demonstrated the practical value of expert systems in business. 

% Building on symbolic logic, expert systems were developed to emulate domain experts' reasoning processes and solve specialized problems \cite{waterman1985guide, jackson1986introduction, giarratano1998expert, mcdermott1982r1, barker1989expert}. 
% These systems use rules to mimic human expertise, providing solutions in fields such as medicine, chemistry, and engineering.

% MYCIN, created in the 1970s, is one of the earliest and most well-known expert systems. 
% It diagnoses blood infections by analyzing symptoms and lab results, offering treatment recommendations through a rule-based engine \cite{van1978mycin}. 
% DENDRAL, which analyzes chemical data like mass spectrometry to predict molecular structures \cite{lindsay1993dendral}, pioneered expert system applications in scientific research, showing how AI could aid in data analysis. 
% COPES, a medical expert system, assists doctors with pathological diagnosis by applying medical knowledge and expert rules \cite{fieschi2013artificial}. 
% By mimicking specialists' reasoning, COPES aimed to improve diagnostic accuracy and support clinical decision-making.




% \subsection{Statistical and Probabilistic Methods}\label{statistical_prob}


% Statistical and probabilistic methods, such as Bayesian inference, Markov processes, and Monte Carlo Tree Search (MCTS), use probabilistic models and manually designed features to handle uncertainty, effectively addressing scenarios with incomplete or noisy data \cite{pearl2014probabilistic}.


% \subsubsection{Bayesian Inference}\label{bayesian}

% Bayesian inference is crucial for modeling dependencies between random variables and their evolution over time. Using frameworks like Bayesian Networks (BNs) and Dynamic Bayesian Networks (DBNs), it allows systems to reason probabilistically about uncertainty. 
% BNs, as directed acyclic graphs, represent conditional dependencies and are used for static causal reasoning \cite{friedman1997bayesian, heckerman1998tutorial}. For instance, in medical diagnosis \cite{weiss1978model}, BNs estimate disease probabilities based on symptoms, while in fault diagnosis, they identify root causes through conditional probabilities. 
% DBNs extend BNs to temporal contexts, making them ideal for modeling time series \cite{murphy2002dynamic}. 
% By linking BNs across time steps, DBNs capture dynamic behaviors and are widely applied in fields like financial forecasting, robotic navigation, and dynamic environment perception.





% \subsubsection{Markov Processes}\label{markov}

% Markov processes model the dynamic evolution of a system's state over time, characterized by the ``Markov property'', where future states depend on the current state, independent of past states \cite{dynkin1965markov}. These include Markov Chains (MCs), Hidden Markov Models (HMMs), and Markov Decision Processes (MDPs), each suited to different applications.

% A MC is a discrete-state model that uses a transition probability matrix to describe state changes \cite{billingsley1961statistical, norris1998markov}, with transitions depending only on the current state:
% \begin{equation}
% \footnotesize
% P(s_{t+1} = j | s_t = i,  s_{t-1},  \dots) = P(s_{t+1} = j | s_t = i)
% \end{equation}
% This model is widely used in weather forecasting, where future states depend on current conditions, and in web page navigation, as in Google PageRank.


% HMMs extend MCs by introducing hidden states, which are inferred from observable data \cite{rabiner1986introduction, eddy1996hidden}. Key components include the initial probability distribution, transition probabilities, and emission probabilities. HMMs are widely used in speech recognition to infer hidden phonemes and in bioinformatics to identify patterns in genes \cite{mor2021systematic}.

% MDPs extend MCs by incorporating actions and rewards to model decision-making, aiming to find a policy that maximizes cumulative rewards \cite{puterman1990markov}. MDPs are widely used in robotic navigation and resource allocation, where the goal is to find optimal paths or strategies to maximize rewards.


\subsection{Monte Carlo Tree Search}\label{mcts}

MCTS is a simulation-based search algorithm for decision-making and planning \cite{browne2012survey}. 
It constructs a search tree through four steps: \textit{Selection}, which chooses the child node with the highest priority using the UCB1 formula:
\begin{equation}
\footnotesize
\text{UCB1} = \frac{w_i}{n_i} + c \sqrt{\frac{\ln N}{n_i}} \text{, }
\end{equation}
where $w_i$ is the total reward of node $i$, $n_i$ is its visit count, $N$ is the parent node's visit count, and $c$ balances exploration and exploitation. \textit{Expansion} adds new nodes, \textit{Simulation} performs random rollouts to evaluate them, and \textit{Backpropagation} updates node statistics. 
MCTS has been widely used in tasks such as optimizing strategies in board games like Go \cite{gelly2011monte} and in robotic path planning, where it helps robots navigate dynamic environments effectively \cite{swiechowski2023monte}.



\noindent\textbf{Summary:} 
MCTS has played a crucial role in the development of reasoning LLMs, particularly in Structural Search (Section \ref{structure_search}). 
By simulating potential future reasoning paths and backpropagating estimated rewards, MCTS helps foundational LLMs efficiently identify the most promising, high-reward paths. 
This process mirrors human-like planning, where future consequences of decisions are considered before taking action. 
By dynamically exploring multiple reasoning trajectories, MCTS enables models to avoid getting stuck in suboptimal paths, making it easier to navigate complex decision spaces. 
This integration has significantly enhanced the ability of LLMs to handle intricate and dynamic reasoning problems, such as those requiring long-term planning or multi-step logical inferences. 
It has allowed LLMs to make more strategic and informed decisions, improving their overall performance in tasks that involve nuanced reasoning and strategic exploration.


% Statistical methods, like MCTS, played a crucial role in the success of DNNs, especially in applications like AlphaZero \cite{silver2017mastering}. 
% By simulating future states and iteratively backpropagating rewards, MCTS enables effective decision-making in complex environments, optimizing strategies such as move exploration in Go. 

% Statistical methods improve AI's ability to handle uncertainty and complex, dynamic data by using probability distributions and random sampling for decision-making in uncertain environments. 
% However, their reliance on manually designed features and model structures can limit their effectiveness in high-dimensional or complex scenarios. 

% With the growth of data availability and computational power, DNNs emerged, marking a transformative shift. 
% DNNs enable end-to-end learning, automatic feature extraction, and the discovery of patterns in large-scale data while capturing nonlinear relationships. 
% This transition from statistical learning to DL signifies a major advancement in AI's sophistication and capability.





\begin{figure*}[t]
    \centering
    \includegraphics[width=0.95\linewidth]{images/compare.pdf}
    \caption{A comprehensive comparison of traditional reasoning models and reasoning LLMs. Reasoning LLMs offer significant advantages over traditional models in areas such as training approaches, adaptability and learning, problem-solving strategies, and generality and scalability.}
    \label{fig:compare}
\end{figure*}




\subsection{Reinforcement Learning}\label{rl}

RL is a type of machine learning where an agent learns to make decisions by interacting with an environment and receiving feedback in the form of rewards, aiming to maximize cumulative rewards over time \cite{sutton1998reinforcement}. 
Early breakthroughs in RL, such as Q-learning \cite{watkins1992q} and DQNs \cite{mnih2015human}, revolutionized the field by enabling the handling of complex state spaces using Deep Neural Networks (DNNs) \cite{torrado2018deep}. These methods paved the way for scaling RL to real-world tasks, where traditional tabular approaches fell short. 
The advent of deep RL marked a significant step forward, combining the power of deep learning with RL to process high-dimensional inputs, such as images and unstructured data.

A landmark achievement in deep RL was AlphaGo, which demonstrated RL's potential by defeating a world champion in the complex game of Go through self-play \cite{silver2016mastering}. 
This success highlighted deep RL's ability to thrive in environments with large, continuous action spaces and uncertainty. 
Building on this, AlphaZero advanced the approach by mastering multiple board games—chess, Go, and Shogi—using self-play, MCTS, and DNNs \cite{silver2017mastering}. 
AlphaZero's ability to learn entirely from scratch, without prior human knowledge, showcased RL's power in environments requiring long-term strategy and planning.

AlphaStar further expanded the boundaries of deep RL by excelling in the real-time strategy game StarCraft II. 
Unlike board games, StarCraft II presents dynamic, partially observable environments and demands multi-step, real-time decision-making \cite{vinyals2019grandmaster}. 
AlphaStar's success in this domain demonstrated deep RL's capacity to adapt to complex decision-making scenarios that require both strategic planning and tactical execution. 
These advancements in RL and deep RL have greatly expanded AI's potential, transitioning from well-defined, static environments to dynamic, complex settings that demand continuous learning and adaptation.


\noindent\textbf{Summary:} 
Deep RL has proven highly effective in solving complex decision-making tasks. AlphaGo exemplifies this by learning strategies through self-play and defeating the world champion in Go. 
This self-play concept laid the foundation for Self Improvement technology (Section \ref{self-improve}) in reasoning LLMs, both relying on continuous feedback and adjustments to optimize strategies.

In RL, reward shaping has been crucial, especially for multi-step reasoning tasks \cite{ng1999policy}. 
By adjusting the reward signal to provide more granular feedback during intermediate steps, it helps agents navigate complex decision-making paths. 
This concept inspired the development of Reward Modeling (Section \ref{prm}), particularly the process reward model, in reasoning LLMs. 
This model offers step-by-step supervision to identify and correct errors in the reasoning process. 
By mimicking human reasoning, the process reward model ensures more robust and interpretable results, especially in tasks like mathematical problem-solving and code generation, where step-by-step evaluation is critical.

Moreover, RL itself is a powerful tool for reasoning LLMs (Section \ref{rl_supervise}). 
With a reward mechanism, RL guides foundational LLMs to find optimal solutions, especially in dynamic reasoning problems. 
Its simplicity and efficiency make RL invaluable for training and optimizing reasoning LLMs, enhancing the intelligence and self-evolution of AI models. 
The integration of RL has led to significant advancements in reasoning LLMs, as demonstrated by DeepSeek-R1 \cite{Deepseek-R1}, offering more flexible and efficient solutions.


% DNNs have made significant strides across several domains, particularly in CV, NLP, and RL.

% \subsubsection{CV}

% In CV, CNNs \cite{lecun1995convolutional} enabled automatic feature extraction, advancing tasks like image classification \cite{deng2009imagenet}, object detection \cite{zhao2019object}, and segmentation \cite{minaee2021image}. 
% ResNets \cite{he2016deep} addressed the vanishing gradient problem in DNNs, improving image processing accuracy. The development of GANs \cite{goodfellow2020generative} has spurred innovation in image generation and style transfer.

% \subsubsection{NLP}

% In NLP, word embedding techniques, such as Word2Vec \cite{mikolov2013distributed} and GloVe \cite{pennington2014glove}, improved semantic understanding by mapping words to low-dimensional spaces. Seq2Seq models \cite{sutskever2014sequence} and LSTMs \cite{hochreiter1997long} advanced machine translation and text generation. 
% Transformer \cite{vaswani2017attention} significantly boosted performance in NLP tasks, especially question answering and dialogue generation.


% \subsubsection{RL}







% \subsection{Foundational LLMs}\label{f_llms}




















% \input{sections/llm_srs_features}





\section{Blueprinting Reasoning LLMs}\label{replication}

In this section, we first analyze the features of reasoning LLMs from both output behavior and training dynamics perspectives. We then provide a detailed overview of the core methods that enable their advanced reasoning capabilities. Finally, we summarize the evolution of reasoning LLMs. 
A comprehensive comparison of traditional reasoning models and reasoning LLMs is shown in Figure \ref{fig:compare}.





\subsection{Analysis of the Features of Reasoning LLMs}\label{o1_features}


\subsubsection{Output Behaviour Perspective}\label{output_behaviour}

\textbf{Explore and Planning Structure:} Recent empirical studies have revealed that reasoning LLMs demonstrate a strong tendency for exploratory behavior in their output structures, especially when compared to models such as WizardMath \cite{wizardmath} and DeepSeekMath \cite{deepseekmath}, which primarily rely on conventional CoT reasoning approaches. 
This exploratory behavior is evident in the formulation of novel hypotheses and the pursuit of alternative solution paths. 
Research by \cite{TowardsSystem2ReasoninLLM} suggests that slow-thinking models engage in a latent generative process, particularly noticeable during the prediction of subsequent tokens. 
This claim is supported by \cite{Deepseek-R1}, which observes that similar behaviors naturally arise during RL scale training. 
Furthermore, the Quiet-STaR framework \cite{QuietStar} introduces an auxiliary pre-training phase focused on next-token prediction, highlighting the critical role of internal deliberation and exploratory mechanisms prior to content generation. 
Collectively, these findings underscore the complex and dynamic nature of reasoning processes in advanced LLMs, emphasizing the interaction between exploration and structured reasoning within their operational frameworks.

\noindent\textbf{Verification and Check Structure:} Analysis of OpenAI's o1 \cite{openai_o1} and o3 \cite{o3-mini} models indicates that their reasoning frameworks incorporate both macro-level actions for long-term strategic planning and micro-level actions, including ``\textit{Wait}'', ``\textit{Hold on}'', ``\textit{Alternatively}'', and ``\textit{Let’s pause}''. 
These micro actions facilitate meticulous verification and iterative checking processes, ensuring precision in task execution. 
Such a dual-layered approach underscores the models' capacity to balance overarching goals with granular, detail-oriented operations, thereby enhancing their overall functionality and reliability. 
To emulate this characteristic, Marco-o1 \cite{Marco_o1}, during the MCTS process for constructing Long-CoT, assigns each tree node the state of ``\textit{Wait! Maybe I made some mistakes! I need to rethink from scratch}'', thereby facilitating the reflective nature of Long-CoT. 
Huatuo-o1 \cite{Huatuo-o1} employs a multi-agent framework to address the issue of incorrect CoT generation during validation. 
This is achieved by incorporating a prompt with ``\textit{Backtracking}'' and ``\textit{Correction}'' functionalities,  which enables the correction process.

\noindent\textbf{Longer Inference Length \& Time:} 
Recent research \cite{TowardsSystem2ReasoninLLM, o1_Journey_Part1, o1_Journey_Part2, huang2025o1, reflection_window} indicates that reasoning LLMs often generate outputs exceeding 2000 tokens to tackle complex problems in coding and mathematics. 
However, this extended output length can sometimes lead to overthinking, where the model spends excessive time on a problem without necessarily improving the solution. 
Studies \cite{TowardsSystem2ReasoninLLM} highlight that while autoregressive generation and Classic CoT can effectively solve simpler problems, they struggle with more complex tasks. 
Research \cite{VisualSlowAgent, SlowPerception} shows that in multimodal domains, many problems demand careful observation, comparison, and deliberation. 
Additionally, Search-o1 \cite{li2025search} suggests that slow-thinking mechanisms are particularly beneficial in areas requiring external knowledge or where potential knowledge conflicts arise. 
In medical scenarios, complex problems, such as those requiring test-time scaling techniques, demonstrate significant improvements \cite{huang2025o1}.



\noindent\textbf{Overly Cautious \& Simple Problem Trap:} Currently, reasoning LLMs have demonstrated strong performance in domains such as competitive-level mathematics \cite{Deepseek-R1, qwq-32b-preview, RedStar, sky_t1_2025}, complex coding \cite{o1_coder}, medical question answering \cite{Huatuo-o1, huang2025o1}, and multilingual translation \cite{Marco_o1, DRT-o1}. 
These scenarios require the model to perform fine-grained analysis of the problem and execute careful logical reasoning based on the given conditions. 
Interestingly, even for straightforward problems like ``\textit{2+3=?}'', reasoning LLMs can exhibit overconfidence or uncertainty. 
Recent research \cite{Tecent_2_plus_3} notes that o1-like models tend to generate multiple solution rounds for easier math problems, often exploring unnecessary paths. 
This behavior contrasts with the lack of diverse exploratory actions for simpler questions, indicating a potential inefficiency in the model's reasoning process.


\subsubsection{Training Dynamic Perspective}\label{dynamic_perspective}

\textbf{Amazing Data Efficiency:} Unlike traditional approaches that focus on expanding instruction sets with uniformly distributed difficulty levels, Studies \cite{huang2025o1, RedStar} suggest that constructing Slow-thinking CoT datasets with a focus on hard samples leads to better generalization in fields like medicine and mathematics. 
This approach diverges from the conventional practice of collecting diverse and evenly distributed instruction datasets.



\noindent\textbf{Sparse Training Method:} 
Contrary to conventional wisdom, the development of effective reasoning LLMs does not require extensive datasets or dense reward signals. 
For example, STILL2 \cite{o1_Journey_Part2} demonstrated impressive performance using only 5,000 distilled samples, while Sky-T1 \cite{sky_t1_2025} achieved performance parity with QwQ \cite{qwq-32b-preview} using just 17,000 LongCoT samples. 
Similarly, RedStar \cite{RedStar} achieved exceptional results across both textual and multimodal tasks with only 4,000 core LongCoT samples. 
In comparison to simple CoT, Slow-thinking Supervised Fine-Tuning (SFT) data exhibits remarkable sample efficiency, often delivering comparable results with just 1/100th of the sample size. 
Additionally, research \cite{simplerl_reason_blob} emphasizes the significant training potential of online RL scaling algorithms, suggesting that non-dense RL supervision and even rule-based reward structures are sufficient for achieving high performance.

\begin{figure}[t]
    \centering
    \includegraphics[width=0.95\linewidth]{images/techs.pdf}
    \caption{The core methods enabling reasoning LLMs.}
    \label{fig:techs}
\end{figure}



\noindent\textbf{Parameter Characteristic:} 
Training LLMs for slow-thinking, as characterized by the LongCoT approach, results in relatively uniform gradient norms across different layers. 
In contrast, fast-thinking, exemplified by the simplified CoT method, generates larger gradient magnitudes in the earlier layers, along with significant variability in gradient norms across layers. 
Empirical evidence suggests that larger models, particularly those exceeding 30 billion parameters, are more compatible with reasoning LLMs training due to their enhanced capacity for complex reasoning. 
Additionally, experiments conducted by RedStar \cite{RedStar} show that the benefits of data scaling vary across model sizes, with scaling effects being more pronounced and effective in larger models. 
This finding is supported by Deepseek-R1's research \cite{Deepseek-R1}, which demonstrates that a 670-billion-parameter model achieves performance metrics closely approximating those of the o1 benchmark, highlighting the scalability advantages of larger architectures in advanced reasoning tasks.








\subsection{Core Method}\label{foundations}

In this section, we provide an overview of the core methods that drive the advanced reasoning capabilities of reasoning LLMs, as shown in Figure \ref{fig:techs}. These include Structure Search, Reward Modeling, Self Improvement, Macro Action, and Reinforcement Fine-Tuning. 
We also highlight representative reasoning LLMs for each method.







\begin{table*}[tbp]
\centering
\renewcommand\arraystretch{1.2}
\caption{Summary of Structure Search method based on the definition of actions and rewards.}
\resizebox{0.98\linewidth}{!}{
\begin{tabular}{l|lp{26em}ll}
\toprule[1.2pt]
 & {{\textbf{Category}}} & \textbf{Reasoning LLMs}  & \textbf{Characteristic}\\
\hline
\multirow{4}{*}{\rotatebox{90}{{\textbf{Actions}}}} & \cellcolor[rgb]{ .949,  .949,  .949} {Reasoning Steps as Nodes} & \cellcolor[rgb]{ .949,  .949,  .949} RAP \cite{hao2023reasoning}, ORM \cite{wan2024alphazero}, Forest-of-Thought \cite{DBLP:journals/corr/abs-2412-09078}  & \cellcolor[rgb]{ .949,  .949,  .949} Actions represent intermediate reasoning steps. \\ 

& {Token-level Decisions}     &  CodeTree \cite{DBLP:journals/corr/abs-2411-04329}, SPaR \cite{DBLP:journals/corr/abs-2412-11605}, TreeBoN \cite{DBLP:journals/corr/abs-2410-16033}  &  Actions involve generating tokens.              \\

& \cellcolor[rgb]{ .949,  .949,  .949} {Task-specific Structures}      & \cellcolor[rgb]{ .949,  .949,  .949} CWM \cite{DBLP:journals/corr/abs-2405-15383}, LLM-MCTS \cite{DBLP:conf/nips/ZhaoLH23} & \cellcolor[rgb]{ .949,  .949,  .949} Actions are domain-specific.                 \\

&  {Correction and Exploration}  &  RethinkMCTS \cite{DBLP:journals/corr/abs-2409-09584}, MCTSr \cite{DBLP:journals/corr/abs-2406-07394} & Actions focus on revisiting and refining previous steps.
\\ %Actions emphasize revisiting, refining, or backtracking to improve previous reasoning steps.

\hline

\multirow{4}{*}{\rotatebox{90}{{\textbf{Rewards}}}} &\cellcolor[rgb]{ .949,  .949,  .949}  {Outcome-based Rewards}     & \cellcolor[rgb]{ .949,  .949,  .949} MC-NEST \cite{DBLP:journals/corr/abs-2411-15645}  & \cellcolor[rgb]{ .949,  .949,  .949} Correctness or validity of the final outcome.  \\ 

&  {Stepwise Evaluations}          & RAP \cite{hao2023reasoning}, SRA-MCTS \cite{DBLP:journals/corr/abs-2411-11053}    &  Rewards are assigned at intermediate steps. \\ 

& \cellcolor[rgb]{ .949,  .949,  .949} {Self-evaluation Mechanisms}              & \cellcolor[rgb]{ .949,  .949,  .949} SPaR \cite{DBLP:journals/corr/abs-2412-11605}, TreeBoN \cite{DBLP:journals/corr/abs-2410-16033}, MindStar \cite{DBLP:journals/corr/abs-2405-16265} & \cellcolor[rgb]{ .949,  .949,  .949} Rewards rely on the model's own confidence.    \\ 

&  {Domain-specific Criteria}          &  LLM-MCTS \cite{DBLP:conf/nips/ZhaoLH23}, SR-MCTS \cite{DBLP:journals/corr/abs-2411-04459}  &  Rewards are tailored to specific tasks. \\ 

& \cellcolor[rgb]{ .949,  .949,  .949} {Iterative Preference Learning}     &\cellcolor[rgb]{ .949,  .949,  .949} LLaMA-Berry \cite{DBLP:journals/corr/abs-2410-02884}, Marco-o1 \cite{Marco_o1}, ReST-MCTS* \cite{DBLP:journals/corr/abs-2406-03816}   & \cellcolor[rgb]{ .949,  .949,  .949} Rewards derive from comparing multiple solutions.\\

\bottomrule[1.2pt]
\end{tabular}
}
\label{tab:structure_search}
\end{table*}



% \begin{table*}[tbp]
% \centering
% \caption{Categories of Structure Search technology based on the definition of actions and rewards.}
% \begin{small}
% \resizebox{\textwidth}{!}{%
% \begin{tabular}{@{}lll@{}}
% \toprule
% \rowcolor[HTML]{FFFFFF} 
% \textbf{Method} &
%   \multicolumn{1}{c}{\cellcolor[HTML]{FFFFFF}\textbf{Category}} &
%   \multicolumn{1}{c}{\cellcolor[HTML]{FFFFFF}\textbf{Characteristic}} \\ \midrule
% \rowcolor[HTML]{ECF4FF} 
% \multicolumn{3}{c}{\cellcolor[HTML]{ECF4FF}\textbf{Actions-based}}                                                                              \\ \midrule
% \rowcolor[HTML]{FFFFFF} 
% RAP \cite{hao2023reasoning}               & Reasoning Steps as Nodes      & \cellcolor[HTML]{FFFFFF}                                                                    \\
% \rowcolor[HTML]{FFFFFF} 
% ORM \cite{wan2024alphazero}               & Reasoning Steps as Nodes      & \cellcolor[HTML]{FFFFFF}                                                                    \\
% \rowcolor[HTML]{FFFFFF} 
% Forest-of-Thought \cite{DBLP:journals/corr/abs-2412-09078} & Reasoning Steps as Nodes      & \multirow{-3}{*}{\cellcolor[HTML]{FFFFFF}Actions represent intermediate reasoning steps.}   \\ \midrule
% \rowcolor[HTML]{EFEFEF} 
% CodeTree \cite{DBLP:journals/corr/abs-2411-04329}         & Token-level Decisions         & \cellcolor[HTML]{EFEFEF}                                                                    \\
% \rowcolor[HTML]{EFEFEF} 
% SPaR \cite{DBLP:journals/corr/abs-2412-11605}              & Token-level Decisions         & \cellcolor[HTML]{EFEFEF}                                                                    \\
% \rowcolor[HTML]{EFEFEF} 
% TreeBoN \cite{DBLP:journals/corr/abs-2410-16033}           & Token-level Decisions         & \multirow{-3}{*}{\cellcolor[HTML]{EFEFEF}Actions involve generating tokens.}                \\ \midrule
% \rowcolor[HTML]{FFFFFF} 
% CWM \cite{DBLP:journals/corr/abs-2405-15383}               & Task-specific Structures      & \cellcolor[HTML]{FFFFFF}                                                                    \\
% \rowcolor[HTML]{FFFFFF} 
% LLM-MCTS \cite{DBLP:conf/nips/ZhaoLH23}         & Task-specific Structures      & \multirow{-2}{*}{\cellcolor[HTML]{FFFFFF}Actions are domain-specific.}                      \\ \midrule
% \rowcolor[HTML]{EFEFEF} 
% RethinkMCTS \cite{DBLP:journals/corr/abs-2409-09584}      & Correction and Exploration    & \cellcolor[HTML]{EFEFEF}                                                                    \\
% \rowcolor[HTML]{EFEFEF} 
% MCTSr \cite{DBLP:journals/corr/abs-2406-07394} &
%   Correction and Exploration &
%   \multirow{-2}{*}{\cellcolor[HTML]{EFEFEF}Actions emphasize revisiting, refining, or backtracking to improve previous reasoning steps.} \\ \midrule
% \rowcolor[HTML]{ECF4FF} 
% \multicolumn{3}{c}{\cellcolor[HTML]{ECF4FF}\textbf{Rewards-based}}                                                                              \\ \midrule
% \rowcolor[HTML]{FFFFFF} 
% MC-NEST \cite{rabby2024mcnestenhancingmathematical}          & Outcome-based Rewards         & Correctness or validity the final outcome.                                                  \\ \midrule
% \rowcolor[HTML]{EFEFEF} 
% RAP \cite{hao2023reasoning}               & Stepwise Evaluations          & \cellcolor[HTML]{EFEFEF}                                                                    \\
% \rowcolor[HTML]{EFEFEF} 
% SRA-MCTS \cite{DBLP:journals/corr/abs-2411-11053}         & Stepwise Evaluations          & \multirow{-2}{*}{\cellcolor[HTML]{EFEFEF}Rewards are assigned at intermediate steps.}       \\ \midrule
% \rowcolor[HTML]{FFFFFF} 
% SPaR \cite{DBLP:journals/corr/abs-2412-11605}             & Self-evaluation Mechanisms    & \cellcolor[HTML]{FFFFFF}                                                                    \\
% \rowcolor[HTML]{FFFFFF} 
% TreeBoN \cite{DBLP:journals/corr/abs-2410-16033}          & Self-evaluation Mechanisms    & \cellcolor[HTML]{FFFFFF}                                                                    \\
% \rowcolor[HTML]{FFFFFF} 
% MindStar \cite{DBLP:journals/corr/abs-2405-16265}          & Self-evaluation Mechanisms    & \multirow{-3}{*}{\cellcolor[HTML]{FFFFFF}Rewards rely on the model's own confidence.}       \\ \midrule
% \rowcolor[HTML]{EFEFEF} 
% LLM-MCTS \cite{DBLP:conf/nips/ZhaoLH23}         & Domain-specific Criteria      & \cellcolor[HTML]{EFEFEF}                                                                    \\
% \rowcolor[HTML]{EFEFEF} 
% SR-MCTS \cite{DBLP:journals/corr/abs-2411-04459}          & Domain-specific Criteria      & \multirow{-2}{*}{\cellcolor[HTML]{EFEFEF}Rewards are tailored to specific tasks.}           \\ \midrule
% \rowcolor[HTML]{FFFFFF} 
% LLaMA-Berry \cite{DBLP:journals/corr/abs-2410-02884}       & Iterative Preference Learning & \cellcolor[HTML]{FFFFFF}                                                                    \\
% \rowcolor[HTML]{FFFFFF} 
% Marco-o1 \cite{Marco_o1}         & Iterative Preference Learning & \cellcolor[HTML]{FFFFFF}                                                                    \\
% \rowcolor[HTML]{FFFFFF} 
% ReST-MCTS* \cite{DBLP:journals/corr/abs-2406-03816}        & Iterative Preference Learning & \multirow{-3}{*}{\cellcolor[HTML]{FFFFFF}Rewards derive from comparing multiple solutions.} \\ \bottomrule
% \end{tabular}%
% }
% \label{tab:structure_search}
% \end{small}
% \end{table*}




\subsubsection{Structure Search}\label{structure_search}
 
Reasoning LLMs aim to achieve high accuracy and depth in solving complex problems by emulating the deliberate nature of human reasoning. 
However, despite recent advancements, current foundational LLMs face inherent limitations when addressing intricate reasoning tasks. 
These limitations arise from their lack of an internal world model to simulate environmental states, their inability to predict the long-term outcomes of reasoning paths, and their failure to iteratively refine reasoning steps based on future states or rewards \cite{huang2023towards}. 
As a result, these shortcomings hinder foundational LLMs from effectively balancing exploration and exploitation in vast reasoning spaces, creating challenges in tasks that require multi-step reasoning, such as complex mathematics, logical inference, or strategic decision-making \cite{DBLP:journals/corr/abs-2405-00451}.

MCTS, a powerful search and optimization algorithm, effectively addresses these challenges by providing a structured framework to explore and evaluate reasoning paths systematically. 
It operates by constructing a reasoning tree, where each node represents a reasoning state, and actions expand the tree by considering potential next steps. 
Through the simulation of future states and the iterative backpropagation of estimated rewards, MCTS allows foundational LLMs to efficiently identify high-reward reasoning paths, mirroring human planning processes. 
This approach aligns with the core principles of reasoning LLMs, where thorough analysis and deliberate exploration are essential for generating well-reasoned outputs. 
Recent methods, such as RAP \cite{hao2023reasoning}, enhance foundational LLMs by integrating MCTS with a world model, enabling the system to iteratively refine intermediate reasoning steps and improve future predictions. 
Similarly, Forest-of-Thought \cite{DBLP:journals/corr/abs-2412-09078} utilizes MCTS to dynamically explore multiple reasoning trajectories, revisiting flawed paths and refining outcomes.





The application of MCTS in reasoning tasks extends beyond traditional problem-solving to highly specialized domains. 
For example, frameworks like SRA-MCTS \cite{DBLP:journals/corr/abs-2411-11053} and MC-NEST \cite{DBLP:journals/corr/abs-2411-15645} showcase the utility of MCTS in tackling technical challenges such as code generation and mathematical reasoning, where intermediate steps are iteratively evaluated and refined. 
In fields like instructional alignment, frameworks such as SPaR \cite{DBLP:journals/corr/abs-2412-11605} and Marco-o1 \cite{Marco_o1} leverage MCTS to refine responses and align reasoning trajectories with human preferences or desired outcomes. 
Additionally, task-specific implementations like HuatuoGPT-o1 \cite{Huatuo-o1} underscore MCTS's crucial role in navigating highly specialized domains, such as medical reasoning, where accuracy and robustness are paramount.

MCTS also enables models to go beyond single-pass reasoning methods, such as CoT or Tree-of-Thought, by incorporating mechanisms to revisit, critique, and refine reasoning steps dynamically \cite{DBLP:journals/corr/abs-2407-01476, DBLP:journals/corr/abs-2409-09584}. 
This iterative capability is essential for tackling tasks with vast decision spaces or those requiring long-term planning, where earlier decisions can significantly impact final outcomes. 
By allowing LLMs to simulate, evaluate, and refine multiple reasoning paths, MCTS introduces a level of adaptability and strategic exploration that traditional approaches lack. 
As shown by AlphaZero-like tree-search \cite{wan2024alphazero} and Search-o1 \cite{li2025search}, MCTS enables reasoning LLMs to not only achieve better performance on specific tasks but also exhibit enhanced generalization capabilities across diverse domains.





The integration of MCTS into LLMs depends on defining actions and rewards to guide reasoning path exploration and assess quality. As shown in Table \ref{tab:structure_search}, we classify the actions in prior work into four categories: 
\begin{enumerate}[itemindent=0em]
\item \textbf{Reasoning Steps as Nodes:} Actions represent intermediate reasoning steps or decisions, such as selecting rules, applying transformations, or generating sub-questions \cite{hao2023reasoning, DBLP:journals/corr/abs-2405-00451, wan2024alphazero, DBLP:journals/corr/abs-2412-09078}.

\item \textbf{Token-level Decisions:} Actions involve generating tokens or sequences (\emph{e.g.}, the next word, phrase, or code snippet) \cite{liu2024dontthrowawayvalue, DBLP:journals/corr/abs-2411-04329, DBLP:journals/corr/abs-2412-11605, DBLP:journals/corr/abs-2410-16033}.

\item \textbf{Task-specific Structures:} Actions are domain-specific, such as moving blocks in blocksworld, constructing geometry in geometry problem-solving, or modifying workflows in task planning \cite{DBLP:journals/corr/abs-2412-10673, DBLP:journals/corr/abs-2405-15383, DBLP:conf/nips/ZhaoLH23}.

\item \textbf{Self-correction and Exploration:} Actions focus on revisiting, refining, or backtracking to improve previous reasoning steps \cite{jiang2024intrinsicselfcorrectionenhancementmonte, DBLP:journals/corr/abs-2409-09584, DBLP:journals/corr/abs-2406-07394}.
\end{enumerate}




Additionally, as illustrated in Table \ref{tab:structure_search}, we classify the reward design into five categories:
\begin{enumerate}[itemindent=0em]
\item \textbf{Outcome-based Rewards:} Rewards focus on the correctness or validity of the final outcome or solution, including the validation of reasoning paths or task success \cite{DBLP:journals/corr/abs-2405-00451, DBLP:journals/corr/abs-2411-15645, DBLP:journals/corr/abs-2412-10673}.

\item \textbf{Stepwise Evaluations:} Rewards are assigned at intermediate steps based on the quality of each step or its contribution toward the final outcome \cite{hao2023reasoning, DBLP:journals/corr/abs-2411-11053, wan2024alphazero}.

\item \textbf{Self-evaluation Mechanisms:} Rewards rely on the model’s own confidence or self-assessment (\emph{e.g.}, likelihood, next-word probability, or confidence scores) \cite{DBLP:journals/corr/abs-2412-11605, DBLP:journals/corr/abs-2410-16033, DBLP:journals/corr/abs-2405-16265}.

\item \textbf{Domain-specific Criteria:} Rewards are tailored to specific tasks, such as symmetry and complexity in geometry or alignment with human preferences in text generation \cite{DBLP:journals/corr/abs-2412-10673, DBLP:conf/nips/ZhaoLH23, DBLP:journals/corr/abs-2411-04459}.

\item \textbf{Iterative Preference Learning:} Rewards are derived from comparing multiple solutions or reasoning paths, guiding learning dynamically \cite{DBLP:journals/corr/abs-2410-02884, Marco_o1, DBLP:journals/corr/abs-2406-03816}.

\end{enumerate}

\noindent\textbf{Summary:} Despite its advantages, structure search-based (\emph{i.e.}, MCTS) reasoning LLMs often suffer from substantial computational overhead due to the large number of simulations required. This makes them less suitable for tasks that demand real-time decision-making or operate under resource constraints \cite{DBLP:journals/corr/abs-2309-03224}. 
Additionally, the effectiveness of MCTS is highly dependent on well-designed reward mechanisms and action definitions, which can vary significantly across different domains, thus posing challenges to its generalizability \cite{DBLP:journals/apin/KemmerlingLS24}.

\begin{table*}[t!]
\centering
\renewcommand\arraystretch{1.2}
\caption{Summary of Reward Modeling method.}
\resizebox{0.98\linewidth}{!}{
\begin{tabular}{lllllll}
\toprule[1.2pt]
\multirow{2}{*}{\textbf{Category}} & \multirow{2}{*}{\textbf{Methods}} & \multirow{2}{*}{\textbf{Data Source}} & \multicolumn{2}{c}{\textbf{Model Refinement}} & \multirow{2}{*}{\textbf{Applications}}      & \multirow{2}{*}{\textbf{Characteristic}}  \\
\cline{4-5}
                          &                          &                                & \textbf{Strategy}              & \textbf{Learning}       &                                   &                                  \\
                          \hline
\multirow{5}{*}{ORM}      & \cellcolor[rgb]{ .949,  .949,  .949}DIVERSE\cite{li2022making}                 & \cellcolor[rgb]{ .949,  .949,  .949}Prompting                      & \cellcolor[rgb]{ .949,  .949,  .949}Fine-tuning           & \cellcolor[rgb]{ .949,  .949,  .949}SFT            & \cellcolor[rgb]{ .949,  .949,  .949}Multiple Reasoning Tasks        & \cellcolor[rgb]{ .949,  .949,  .949}Weighted Voting   Verifier       \\
                          & MATH-SHEPHERD\cite{wang2024math}            & Sampling                       & Feedback-guided       & SFT \& RL      & Math Reasoning                    & Correctness Score Assignment   \\
                          & \cellcolor[rgb]{ .949,  .949,  .949}AutoPSV\cite{lu2024autopsv}                  & \cellcolor[rgb]{ .949,  .949,  .949}Prompting                      & \cellcolor[rgb]{ .949,  .949,  .949}Feedback-guided       & \cellcolor[rgb]{ .949,  .949,  .949}SFT            & \cellcolor[rgb]{ .949,  .949,  .949}Math / Commonsense Reasoning    & \cellcolor[rgb]{ .949,  .949,  .949}Automated Process Supervision  \\
                          & Implicit PRMs\cite{yuan2024free}            & Sampling                       & Fine-tuning           & SFT \& RL      & Math Reasoning                    & Obtaining PRM from ORM         \\
                          & \cellcolor[rgb]{ .949,  .949,  .949}OVM\cite{yu2024ovm}                      & \cellcolor[rgb]{ .949,  .949,  .949}Sampling                       & \cellcolor[rgb]{ .949,  .949,  .949}Feedback-guided       & \cellcolor[rgb]{ .949,  .949,  .949}SFT            & \cellcolor[rgb]{ .949,  .949,  .949}Math Reasoning                    & \cellcolor[rgb]{ .949,  .949,  .949}Guided Decoding                  \\
                          \hline
\multirow{4}{*}{MCTS}     & ReST-MCTS$^{*}$\cite{zhang2024rest}              & Sampling                       & Self-training         & SFT \& RL      & Multiple Reasoning Tasks        & MCTS and Self-training        \\
                          & \cellcolor[rgb]{ .949,  .949,  .949}OmegaPRM\cite{luo2024improve}                 & \cellcolor[rgb]{ .949,  .949,  .949}MCTS with Binary Search      & \cellcolor[rgb]{ .949,  .949,  .949}Feedback-guided       & \cellcolor[rgb]{ .949,  .949,  .949}SFT            & \cellcolor[rgb]{ .949,  .949,  .949}Math Reasoning                    & \cellcolor[rgb]{ .949,  .949,  .949}Divide-and-Conquer MCTS        \\
                          & ReARTeR\cite{sun2025rearter}                  & Sampling                       & Feedback-guided       & SFT \& RL      & QA                                & Retrieval-Augmented   Generation \\
                          & \cellcolor[rgb]{ .949,  .949,  .949}Consensus Filtering\cite{zhang2025lessons}      & \cellcolor[rgb]{ .949,  .949,  .949}MCTS Data   Construction       & \cellcolor[rgb]{ .949,  .949,  .949}Feedback-guided       & \cellcolor[rgb]{ .949,  .949,  .949}SFT            & \cellcolor[rgb]{ .949,  .949,  .949}Math Reasoning                    & \cellcolor[rgb]{ .949,  .949,  .949}Consensus Filtering Mechanism  \\
                          \hline
\multirow{3}{*}{PRM}      & ORPS\cite{yu2024outcome}                     & Sampling                       & Feedback-guided       & SFT            & Code Generation                   & Supervising Outcome Refinement \\
                          & \cellcolor[rgb]{ .949,  .949,  .949}Step-DPO\cite{lai2024step}                 & \cellcolor[rgb]{ .949,  .949,  .949}Sampling                       & \cellcolor[rgb]{ .949,  .949,  .949}Feedback-guided       & \cellcolor[rgb]{ .949,  .949,  .949}SFT \& RL      & \cellcolor[rgb]{ .949,  .949,  .949}Math Reasoning                    & \cellcolor[rgb]{ .949,  .949,  .949}Step-wise Preference Pairs     \\
                          & AdaptiveStep\cite{liu2025adaptivestep}             & Response Dividing              & Feedback-guided       & SFT            & Math Reasoning, Code Generation & Dividing   Reasoning Steps \\
\bottomrule[1.2pt]
\end{tabular}}
\label{table:PRM}
\end{table*}





% \begin{table*}[t!]
% \centering
% \caption{Summary of PRM technology.}
% \bgroup
% \def\arraystretch{1,2}
% \resizebox{0.9\linewidth}{!}{
% \begin{tabular}{llll}
% \toprule[1.2pt]
% \textbf{Category}   & \textbf{Reasoning LLMs}   & \textbf{LLMs Backbone}           & \textbf{Characteristic}\\ 
% \hline
% \multirow{5}{*}{ORM} 
%  &  \cellcolor[rgb]{ .949,  .949,  .949} DIVERSE\cite{li2022making}
% &\cellcolor[rgb]{ .949,  .949,  .949} DeBERTa-v3/Davinci
% & \cellcolor[rgb]{ .949,  .949,  .949} Weighted Voting Verifier\\

% & MATH-SHEPHERD\cite{wang2024math}
% & LLaMA2/LLEMMA/Mistral/DeepSeek         & Correctness Score Assignment\\   

% &\cellcolor[rgb]{ .949,  .949,  .949}AutoPSV\cite{lu2024autopsv}
% &\cellcolor[rgb]{ .949,  .949,  .949}Phi2/Mistral/Mixtral/Qwen
% &\cellcolor[rgb]{ .949,  .949,  .949}Automated Process Supervision\\

% & Implicit PRMs\cite{yuan2024free}
% & Mistral/LLaMA3 
% & Obtaining PRM from ORM\\

% &\cellcolor[rgb]{ .949,  .949,  .949}OVM\cite{yu2024ovm}
% &\cellcolor[rgb]{ .949,  .949,  .949}LLaMA2/Mistral
% &\cellcolor[rgb]{ .949,  .949,  .949}Guided Decoding\\ 
% \hline

% \multirow{3}{*}{MCTS}
% & ReST-MCTS$^{*}$\cite{zhang2024rest}
% & LLaMA3/Mistral/SciGLM
% & Self Training\\

% &\cellcolor[rgb]{ .949,  .949,  .949}OmegaPRM\cite{luo2024improve}
% & \cellcolor[rgb]{ .949,  .949,  .949}Gemini1.5/Gemma2
% & \cellcolor[rgb]{ .949,  .949,  .949}Divide-and-Conquer MCTS\\

% & ReARTeR\cite{sun2025rearter}
% & GPT-4o/LLaMA3
% & Retrieval-augmented Generation\\
% \hline
% \multirow{2}{*}{PRM} 
% &\cellcolor[rgb]{ .949,  .949,  .949}ORPS\cite{yu2024outcome}
% &\cellcolor[rgb]{ .949,  .949,  .949} LLaMA3/DeepSeek-Coder/Qwen-2.5/GPT-4o
% &\cellcolor[rgb]{ .949,  .949,  .949}Supervising Outcome Refinement\\

% & Step-DPO\cite{lai2024step}
% & Qwen/LLaMA3/DeepSeek-Math
% & Step-wise Preference Pairs   \\

% \bottomrule[1.2pt]
% \end{tabular}}
% \label{table:PRM}
% \egroup
% \end{table*}





\subsubsection{Reward Modeling}\label{prm}


Two primary training paradigms are used to tackle multi-step reasoning tasks: outcome supervision and process supervision. 
Outcome supervision emphasizes the correctness of the final answer at a higher level of granularity, and the resulting model is referred to as the Outcome Reward Model (ORM) \cite{cobbe2021training, yu2023outcome}. 
In contrast, process supervision provides step-by-step labels for the solution trajectory, evaluating the quality of each reasoning step. 
The resulting model is known as the Process Reward Model (PRM) \cite{uesato2022solving, lightmanlet, li2023making}. The main distinction between ORM and PRM is illustrated in Figure \ref{fig:prm_vs_orm}.


\begin{figure}[t]
    \centering
    \includegraphics[width=0.95\linewidth]{images/PRM_ORM.pdf}
    \caption{The comparison between ORM and PRM for assessing a complete solution trajectory. ORM only provides a single reward based on the correctness of the final answer, while PRM evaluates the quality of each reasoning step throughout the process.}
    \label{fig:prm_vs_orm}
\end{figure}




PRM offers significant advantages \cite{wu2023fine, wang2024math} in complex reasoning tasks for several key reasons. First, it provides fine-grained, step-wise supervision, allowing for the identification of specific errors within a solution path. This feature is especially valuable for RL and automated error correction. Second, PRM closely mirrors human reasoning behavior, which relies on accurate intermediate steps to reach correct conclusions. Unlike ORM, PRM avoids situations where incorrect reasoning can still lead to a correct final answer, thus ensuring more robust and interpretable reasoning. While PRM has primarily been applied to complex mathematical problems, its benefits have recently driven applications in other fields. For instance, ORPS \cite{yu2024outcome} utilizes PRM to address complex code generation challenges, while Step-DPO \cite{lai2024step} combines process supervision with the Direct Preference Optimization (DPO) algorithm \cite{rafailov2024direct} to improve long-chain mathematical reasoning. A summary of Reward Modeling method is presented in Table \ref{table:PRM}.





\noindent\textbf{Summary:} 
Despite the advantages of PRMs, they present several challenges. 
The primary difficulty is obtaining process supervision-labeled data, which is often both costly and time-consuming. 
To address concerns related to scalability, efficiency, and accuracy, researchers have explored various automated annotation methods. 
For example, MATH-SHEPHERD \cite{wang2024math} utilizes the correctness of the final answer to define the quality of intermediate steps based on their potential to lead to the correct outcome, automating the step-wise data collection process. ReST-MCTS$^{*}$ \cite{zhang2024rest} combines process reward guidance with MCTS to generate higher-quality reasoning traces through extensive rollouts. 
Similarly, OmegaPRM \cite{luo2024improve} employs the MCTS framework while introducing a divide-and-conquer algorithm for automated process supervision data generation. 
Another novel approach involves using ORM to train a PRM. Yuan et al. \cite{yuan2024free} propose training a PRM implicitly by leveraging ORM training on cheaper datasets, under mild reward parameterization assumptions. 
They also provide theoretical guarantees for the performance of this implicit PRM, demonstrating its practicality and cost-effectiveness.

In addition to data collection, PRMs face challenges related to trustworthiness \cite{sun2025rearter}, categorized as follows: 
\begin{enumerate}[itemindent=0em]
\item \textbf{Lack of Explanations:} Current PRMs often generate scores for reasoning steps without sufficient explanations, limiting interpretability and hindering their usefulness in refining reasoning during test-time. 

\item \textbf{Bias in Training Data:} Data collection methods, such as MCTS, tend to introduce distributional biases, assigning disproportionately higher scores to the majority of questions. As a result, PRMs struggle to effectively identify erroneous reasoning steps. 

\item \textbf{Early-Step Bias:} PRMs show lower accuracy in predicting rewards for earlier reasoning steps compared to those closer to the final answer. This issue stems from the increased randomness and uncertainty associated with the initial steps in the reasoning process.

\end{enumerate}



\begin{table*}[t!]
\centering
\renewcommand\arraystretch{1.2}
\caption{Summary of Self Improvement method.}

%We summarize the key features for each work: \textbf{Data Source}(Data exploration methods for training), \textbf{Model Refinement}(includes \textbf{Feedback} generation and learning \textbf{Strategy}) and \textbf{Applications}.

\resizebox{0.98\linewidth}{!}{
\begin{tabular}{llllll}
\toprule[1.2pt]
\multirow{2}{*}{\textbf{Stage}}   & \multirow{2}{*}{\textbf{Methods}}   & \multirow{2}{*}{\textbf{Data Source}}  & \multicolumn{2}{c}{\textbf{Model Refinement}} & \multirow{2}{*}{\textbf{Application}}\\\cline{4-5}
& &  & \multicolumn{1}{l}{\textbf{Feedback}} & \multicolumn{1}{l}{\textbf{Strategy}}                            \\ \hline

\multirow{11}{*}{Training} & \cellcolor[rgb]{ .949,  .949,  .949}STaR \cite{DBLP:conf/nips/ZelikmanWMG22}                   & \cellcolor[rgb]{ .949,  .949,  .949}Few-shot  & \cellcolor[rgb]{ .949,  .949,  .949}Language Model & \cellcolor[rgb]{ .949,  .949,  .949}SFT                       & \cellcolor[rgb]{ .949,  .949,  .949}QA, Arithmetic Reasoning                                                        \\

&  Quiet-STaR\cite{QuietStar} &  Token-level Exploration &Language Model &RL                   &  QA, Arithmetic Reasoning                                                                     \\
& \cellcolor[rgb]{ .949,  .949,  .949}V-STaR \cite{hosseini2024vstartrainingverifiersselftaught} & \cellcolor[rgb]{ .949,  .949,  .949}Sampling & \cellcolor[rgb]{ .949,  .949,  .949}Verifier & \cellcolor[rgb]{ .949,  .949,  .949}SFT                   & \cellcolor[rgb]{ .949,  .949,  .949}Arithmetic Reasoning, Code Generation \\

&  B-STaR \cite{zeng2024bstarmonitoringbalancingexploration}  & Sampling &Reward Model &SFT &  Arithmetic Reasoning, Code Generation    \\

& \cellcolor[rgb]{ .949,  .949,  .949}rStar-Math \cite{guan2025rstarmathsmallllmsmaster}         & \cellcolor[rgb]{ .949,  .949,  .949}MCTS Data Construction & \cellcolor[rgb]{ .949,  .949,  .949}Reward Model & \cellcolor[rgb]{ .949,  .949,  .949}SFT                            & \cellcolor[rgb]{ .949,  .949,  .949}Arithmetic Reasoning \\

&  ReST \cite{DBLP:journals/corr/abs-2308-08998}              & Sampling &Reward Model &RL                                          &Machine Translation      \\

& \cellcolor[rgb]{ .949,  .949,  .949}ReST-EM \cite{singh2024humandatascalingselftraining}       & \cellcolor[rgb]{ .949,  .949,  .949}Sampling & \cellcolor[rgb]{ .949,  .949,  .949}Language Model & \cellcolor[rgb]{ .949,  .949,  .949}EM for RL                                   & \cellcolor[rgb]{ .949,  .949,  .949}Arithmetic Reasoning, Code Generation            \\

&  ReST-MCTS* \cite{zhang2024rest}  & Sampling  &Reward Model   &SFT, RL                    &  Reasoning                                              \\

& \cellcolor[rgb]{ .949,  .949,  .949}ENVISIONS \cite{DBLP:journals/corr/abs-2406-11736}         & \cellcolor[rgb]{ .949,  .949,  .949}Sampling & \cellcolor[rgb]{ .949,  .949,  .949}Environment Guided & \cellcolor[rgb]{ .949,  .949,  .949}SFT                      & \cellcolor[rgb]{ .949,  .949,  .949}Web Agents, Reasoning    \\

&  RISE \cite{qu2024recursiveintrospectionteachinglanguage}   &  Sampling &Reward Function &Weighted SFT                                &  Arithmetic Reasoning   \\

& \cellcolor[rgb]{ .949,  .949,  .949}STIC \cite{deng2024enhancinglargevisionlanguage}           & \cellcolor[rgb]{ .949,  .949,  .949}Few-shot & \cellcolor[rgb]{ .949,  .949,  .949}Language Model & \cellcolor[rgb]{ .949,  .949,  .949}SFT                       & \cellcolor[rgb]{ .949,  .949,  .949}Vision Language Model Tasks            \\ 

&  SIRLC \cite{DBLP:conf/iclr/PangWLC0Z024}   &Question Answeing  &Language Model &RL    &Reasoning, Translation, Summary   \\

& \cellcolor[rgb]{ .949,  .949,  .949}AlpacaFarm \cite{DBLP:conf/nips/DuboisLTZGBGLH23}           & \cellcolor[rgb]{ .949,  .949,  .949}Existing Data & \cellcolor[rgb]{ .949,  .949,  .949}Language Model & \cellcolor[rgb]{ .949,  .949,  .949}SFT                       & \cellcolor[rgb]{ .949,  .949,  .949}None (Intrinsic Evaluation)            \\ 

\hline

\multirow{4}{*}{Inference} & Self-Refine \cite{DBLP:conf/nips/MadaanTGHGW0DPY23}   &Independent of Training Data     & Language Model        & Few-shot Demonstration                        & Code Generation, Sentiment Reversal, Acronym Generation                     \\

& \cellcolor[rgb]{ .949,  .949,  .949}Self-Check \cite{DBLP:conf/iclr/MiaoTR24}       & \cellcolor[rgb]{ .949,  .949,  .949}Independent of Training Data            & \cellcolor[rgb]{ .949,  .949,  .949}Language Model           & \cellcolor[rgb]{ .949,  .949,  .949}Step Check               & \cellcolor[rgb]{ .949,  .949,  .949}QA, Arithmetic Reasoning   \\

&  CRITIC \cite{DBLP:conf/iclr/GouSGSYDC24}               &Independent of Training Data    &Language Model  & External Tools                             &  QA, Arithmetic Reasoning, Detoxification          \\

& \cellcolor[rgb]{ .949,  .949,  .949}ROSE\cite{zhong2024rosedoesntthatboosting}              & \cellcolor[rgb]{ .949,  .949,  .949}Independent of Training Data & \cellcolor[rgb]{ .949,  .949,  .949}Language Model                    & \cellcolor[rgb]{ .949,  .949,  .949}Distributed Prompt                     & \cellcolor[rgb]{ .949,  .949,  .949}Safety, Knowledge    \\ 

&  Self-Verification \cite{DBLP:conf/emnlp/WengZX0HLSLZ23}               &Independent of Training Data    &Language Model  & Re-Ranking                             & Arithmetic Reasoning          \\

& \cellcolor[rgb]{ .949,  .949,  .949}SelfEval-Decoding \cite{xie2023selfevaluationguidedbeamsearch}              & \cellcolor[rgb]{ .949,  .949,  .949}Independent of Training Data & \cellcolor[rgb]{ .949,  .949,  .949}Language Model                    & \cellcolor[rgb]{ .949,  .949,  .949}Beam Search                     & \cellcolor[rgb]{ .949,  .949,  .949}Aritnmetic/Symbolic Reasoning                      \\ 

&  IPS \cite{yao2023finegrainedconversationaldecodingisotropic}               &Independent of Training Data    &Language Model  &Constrained Decoding    &Dialogue \\

& \cellcolor[rgb]{ .949,  .949,  .949}Control-DAG\cite{chen2024controldagconstraineddecodingnonautoregressive}              & \cellcolor[rgb]{ .949,  .949,  .949}Independent of Training Data & \cellcolor[rgb]{ .949,  .949,  .949}Language Model                    & \cellcolor[rgb]{ .949,  .949,  .949}Constrained Decoding                     & \cellcolor[rgb]{ .949,  .949,  .949}Dialogue, Open-domain Generation    \\ 

&  Look-Back \cite{xu2023lookbackdecodingopenendedtext}   &Independent of Training Data    &Language Model  &Contrastive Decoding    &Alleviating Repetitions \\


& \cellcolor[rgb]{ .949,  .949,  .949}LeCo \cite{DBLP:journals/corr/abs-2403-19094}              & \cellcolor[rgb]{ .949,  .949,  .949}Independent of Training Data & \cellcolor[rgb]{ .949,  .949,  .949}Language Model                    & \cellcolor[rgb]{ .949,  .949,  .949}Constrained Decoding                     & \cellcolor[rgb]{ .949,  .949,  .949}QA, Reasoning                       \\ 

\bottomrule[1.2pt]
\end{tabular}
}
\label{table:self-improve}
\end{table*}




\subsubsection{Self Improvement}\label{self-improve}

Reasoning LLMs exemplify a progression from weak to strong supervision, while traditional CoT fine-tuning faces challenges in scaling effectively. 
Self improvement, using the model's exploration capabilities for self-supervision, gradually enhances LLMs performance in tasks such as translation \cite{DBLP:journals/corr/abs-2308-08998}, mathematics \cite{DBLP:conf/nips/ZelikmanWMG22, singh2024humandatascalingselftraining}, and multimodal perception \cite{deng2024enhancinglargevisionlanguage}. 
This approach fosters exploration and application within reasoning LLMs \cite{anthony2017thinkingfastslowdeep, tong2024can, tong2024optimizing, guan2025rstarmathsmallllmsmaster}. 
A summary of Self Improvement method is presented in Table \ref{table:self-improve}.








Training-based self improvement in LLMs can be categorized based on exploration and improvement strategies. 
The exploration phase focuses on data collection to facilitate subsequent training improvements, with notable variations in approach. 
STaR \cite{DBLP:conf/nips/ZelikmanWMG22} uses few-shot examples for data gathering, while ReST \cite{DBLP:journals/corr/abs-2308-08998}, ReST-EM \cite{singh2024humandatascalingselftraining}, and ENVISIONS \cite{DBLP:journals/corr/abs-2406-11736} rely on multiple samplings of complete trajectories. 
Quiet-STaR \cite{QuietStar} explores at the token level, introducing concepts like meta-tokens and non-myopic loss to enhance supervision. 
Additionally, ReST-MCTS* \cite{zhang2024rest} and rStar-Math \cite{guan2025rstarmathsmallllmsmaster} generate training data through MCTS.




Improvement strategies also exhibit significant diversity. 
For instance, STaR and its derivatives, such as V-STaR \cite{hosseini2024vstartrainingverifierselftaught} and B-STaR \cite{zeng2024bstarmonitoringbalancingexploration}, combine filtering with SFT. 
ReST and its variants typically introduce innovative reward calculation methods to enhance RL training for policy models. 
RISE \cite{qu2024recursiveintrospectionteachinglanguage} incorporates external feedback, recording rewards and refining responses through distillation during the improvement process. 
Notably, rStar-Math \cite{guan2025rstarmathsmallllmsmaster} demonstrates that small models have achieved \textit{System 2} reflective capabilities through self-evolving training approaches.


Test-time self improvement leverages the consistency of a model's internal knowledge to correct hallucinations during inference. 
These approaches can be categorized into three main types: methods that refine answers using prompts \cite{DBLP:conf/nips/MadaanTGHGW0DPY23, DBLP:conf/iclr/MiaoTR24}, approaches that utilize external tools \cite{DBLP:conf/iclr/GouSGSYDC24}, and techniques that leverage logits without the need for external tools or prompts \cite{DBLP:journals/corr/abs-2403-19094, xu2023lookbackdecodingopenendedtext}.




\subsubsection{Macro Action}\label{macro_action}

Recent advancements in LLMs have driven progress in emulating human-like \textit{System 2} cognitive processes via sophisticated thought architectures, often referred to as macro action frameworks. 
These structured reasoning systems go beyond traditional token-level autoregressive generation by introducing hierarchical cognitive phases, such as strategic planning, introspective verification, and iterative refinement. 
This approach not only enhances the depth of reasoning but also broadens the solution space, enabling more robust and diverse problem-solving pathways. A summary of Macro Action method is presented in Table \ref{table:macro_action}.




% \begin{table*}[!t]
% \centering
% \caption{Summary of Macro Action method. \textbf{Tag}: \encircle[fill=harvestgold, text=white]{I} = \underline{I}mage, \encircle[fill=lightcoral, text=white]{T} = \underline{T}ext, \encircle[fill=DarkGreen, text=white]{V} = \underline{V}ideo.}

% %We summarize the key features for each work: the \textbf{Modality}, \textbf{Reflection}, \textbf{Action Number}, \textbf{Action Source}: Human-designed tempelete or not, the \textbf{Learning}: ICL, RL or SFT, \textbf{Application \& Benchmark}: Summary, Translation, Math or Code. \textbf{Tag}: \encircle[fill=harvestgold, text=white]{I} = \underline{I}mage, \encircle[fill=lightcoral, text=white]{T} = \underline{T}ext, \encircle[fill=DarkGreen, text=white]{V} = \underline{V}ideo 

% \resizebox{0.98\textwidth}{!}{
%    \begin{tabular}% 
%    {lllllccl}
%     \toprule
%         \multirow{2}{*}{\textbf{Methods}} & \multirow{2}{*}{\textbf{Usage}} & \multicolumn{5}{c}{\textbf{Action Attribute}} & \multirow{2}{*}{\textbf{Representative Action}} \\ \cline{3-7}
%           & & \textbf{Action Source} & \textbf{Action Number} & \textbf{Learning} & \textbf{Reflection} & \textbf{Modality} & \\ \midrule
%         % RLHF
%        \rowcolor[rgb]{ .949,  .949,  .949} Self-Check\cite{miao2023selfcheck} & Verification & Human-Designed & 4 & ICL & \textcolor{green}{\ding{51}} & \encircle[fill=lightcoral, text=white]{T} & Target Extraction, Information Collection, Step Regeneration, Result Comparsion \\
       
%         LeMa\cite{LeMa} & Synthetic Data & Human-Designed & 3 & ICL \& SFT & \textcolor{green}{\ding{51}} & \encircle[fill=lightcoral, text=white]{T} & Incorrect Step Recognition, Explanation, Correct Solution: \\
        
%        \rowcolor[rgb]{ .949,  .949,  .949} REFINER\cite{li2024refiner} & Verification/Exploration & Human-Designed & 2 & ICL \& SFT & \textcolor{green}{\ding{51}} & \encircle[fill=lightcoral, text=white]{T} & Critic, Generate \\
        
%         HiICL-MCTS\cite{HiICL-MCTS} & Exploration & Human-Designed & 5 & ICL & \textcolor{green}{\ding{51}} & \encircle[fill=lightcoral, text=white]{T} & System Analysis, One-Step Thought, Divide and Conquer, ..., Self-Reflection and Refinement  \\
        
%          \rowcolor[rgb]{ .949,  .949,  .949} SUPERCORRECT\cite{yang2024supercorrect} & Distill & In-Context Learning & Dynamic & SFT \& RL & \textcolor{red}{\ding{55}} & \encircle[fill=lightcoral, text=white]{T} & -- \\
        
%         ReasonFlux\cite{ReasonFlux}  & Synthetic Data/Exploration & Human-Designed & $\sim$500 & ICL \& SFT \& RL & \textcolor{red}{\ding{55}} &  \encircle[fill=lightcoral, text=white]{T} & -- \\
        
%          \rowcolor[rgb]{ .949,  .949,  .949} rStar\cite{rSTaR} & Exploration & Human-Designed & 5 & ICL \& RL & \textcolor{green}{\ding{51}} & \encircle[fill=lightcoral, text=white]{T} & One-step thought, Propose Next Sub-question \& Answer, ..., Rephrase question \\
         
%         LLaMA-Berry\cite{zhang2024llama} & Exploration & Human-Designed & 2 & ICL \& RL & \textcolor{green}{\ding{51}} & \encircle[fill=lightcoral, text=white]{T} & Reflection, Error Re-correction \\
    
%         \rowcolor[rgb]{ .949,  .949,  .949} Huatuo-o1\cite{Huatuo-o1} & Synthetic Data & Human-Designed & 4 & ICL \& SFT & \textcolor{green}{\ding{51}} & \encircle[fill=lightcoral, text=white]{T} & Backtracking, Exploring New
% Paths, Verification, Correction \\

%         Marco-o1\cite{Marco_o1} & Verification & Human-Designed & 1 & ICL \& SFT & \textcolor{green}{\ding{51}} & \encircle[fill=lightcoral, text=white]{T} & Reflection \\
        
%          \rowcolor[rgb]{ .949,  .949,  .949}BoT\cite{yang2024buffer} & Exploration & In-Context Learning & Dynamic & ICL & \textcolor{red}{\ding{55}} & \encircle[fill=lightcoral, text=white]{T} & Solving Quadratic Equation, Array Sorting, ...,  Search Algorithms) \\
         
%         rStar-Math\cite{guan2025rstarmathsmallllmsmaster} & Exploration &  In-Context Learning & $1$ & ICL \& RL & \textcolor{green}{\ding{51}} & \encircle[fill=lightcoral, text=white]{T} & Python comment \\
        
%          \rowcolor[rgb]{ .949,  .949,  .949}Mulberry\cite{yao2024mulberry} & Synthetic Data & In-Context Learning & 1 & ICL \& SFT & \textcolor{green}{\ding{51}}  & \encircle[fill=lightcoral, text=white]{T} & Reflection \\
         
%         LLaVA-CoT\cite{xu2024llava} & Synthetic Data/Exploration & Human-Designed & 4 & SFT & \textcolor{red}{\ding{55}} & \encircle[fill=harvestgold, text=white]{I} \space \encircle[fill=lightcoral, text=white]{T} & Summary, Caption, Reasoning, Conclusion \\
        
%          \rowcolor[rgb]{ .949,  .949,  .949}LLaMAV-o1\cite{thawakar2025llamav} & Verification/Exploration & Human-Designed & 4173 & Curriculum Learning & \textcolor{green}{\ding{51}} &  \encircle[fill=harvestgold, text=white]{I} \space \encircle[fill=lightcoral, text=white]{T} & Detailed Caption Generation, Logical Reasoning, ... Final Answer Generation \\
         
%         AtomThink\cite{atomthink} & Synthetic Data/Exploration & In-Context Learning & $>$100 & SFT \& RL & \textcolor{green}{\ding{51}} & \encircle[fill=harvestgold, text=white]{I} \space \encircle[fill=lightcoral, text=white]{T} & Variable Definition, Calculations, Graphs Analysis , ...,  Verification\\

%          \rowcolor[rgb]{ .949,  .949,  .949}RedStar\cite{RedStar} & Distill & Human-Designed & 2 & SFT & \textcolor{green}{\ding{51}} & \encircle[fill=harvestgold, text=white]{I} \space \encircle[fill=lightcoral, text=white]{T} & Wait, Alternately \\
        
%         Auto-CoT\cite{zhang2210automatic}  & Exploration & In-Context Learning & 2 & ICL & \textcolor{red}{\ding{55}} & \encircle[fill=lightcoral, text=white]{T} & Question clustering, Demonstration Sampling \\
        
%          \rowcolor[rgb]{ .949,  .949,  .949}PoT\cite{chen2022program} & Verification & In-Context Learning & 1 & ICL & \textcolor{red}{\ding{55}} & \encircle[fill=lightcoral, text=white]{T} & Code language conversion \\
        
%         PAL\cite{gao2023pal} & Verification & In-Context Learning & 1 & ICL & \textcolor{red}{\ding{55}} & \encircle[fill=lightcoral, text=white]{T} & Code language conversion \\
        
%          \rowcolor[rgb]{ .949,  .949,  .949} Decomposed Prompt\cite{decomposed} & Exploration & Human-Designed & 3 & ICL & \textcolor{red}{\ding{55}} & \encircle[fill=lightcoral, text=white]{T} & Peoblem Split, Subproblem Solving, Answer Merge \\
        
%         Least-to-Most\cite{least2most} & Exploration & Human-Designed & 2 & ICL & \textcolor{red}{\ding{55}} & \encircle[fill=lightcoral, text=white]{T} & Problem Decomposition, Subproblem Solving \\
%         \bottomrule
%    \end{tabular}
% }
%     \label{table:macro_action}
% \end{table*}



\begin{table*}[!t]
\centering
\caption{Summary of Macro Action method. \textbf{Tag}: \encircle[fill=harvestgold, text=white]{I} = \underline{I}mage, \encircle[fill=lightcoral, text=white]{T} = \underline{T}ext, \encircle[fill=DarkGreen, text=white]{V} = \underline{V}ideo. Action Category: AD: Analysis and Decomposition, IPR: Information Processing and Reasoning, VC: Verification and Correction, GO: Generation and Optimization, EB: Exploration and Backtracking.}  

\resizebox{0.98\textwidth}{!}{
   \begin{tabular}% 
   {lllllccc}
    \toprule
        \multirow{2}{*}{\textbf{Methods}} & \multirow{2}{*}{\textbf{Usage}} & \multicolumn{5}{c}{\textbf{Action Attribute}} & \multirow{2}{*}{\textbf{Main Action Category}} \\ \cline{3-7}
          & & \textbf{Action Source} & \textbf{Action Number} & \textbf{Learning} & \textbf{Reflection} & \textbf{Modality} & \\ \midrule
        % RLHF
       \rowcolor[rgb]{ .949,  .949,  .949} Self-Check\cite{miao2023selfcheck} & Verification & Human-Designed & 4 & ICL & \textcolor{green}{\ding{51}} & \encircle[fill=lightcoral, text=white]{T} & AD, VC \\
       
        LeMa\cite{LeMa} & Synthetic Data & Human-Designed & 3 & ICL \& SFT & \textcolor{green}{\ding{51}} & \encircle[fill=lightcoral, text=white]{T} & VC, IPR \\
        
       \rowcolor[rgb]{ .949,  .949,  .949} REFINER\cite{li2024refiner} & Verification/Exploration & Human-Designed & 2 & ICL \& SFT & \textcolor{green}{\ding{51}} & \encircle[fill=lightcoral, text=white]{T} & VC, AD \\
        
        HiICL-MCTS\cite{HiICL-MCTS} & Exploration & Human-Designed & 5 & ICL & \textcolor{green}{\ding{51}} & \encircle[fill=lightcoral, text=white]{T} & VC, EB, AD  \\
        
         \rowcolor[rgb]{ .949,  .949,  .949} SUPERCORRECT\cite{yang2024supercorrect} & Distill & In-Context Learning & Dynamic & SFT \& RL & \textcolor{red}{\ding{55}} & \encircle[fill=lightcoral, text=white]{T} & AD, IPR \\
        
        ReasonFlux\cite{ReasonFlux}  & Synthetic Data/Exploration & Human-Designed & $\sim$500 & ICL \& SFT \& RL & \textcolor{red}{\ding{55}} &  \encircle[fill=lightcoral, text=white]{T} & AD, IPR \\
        
         \rowcolor[rgb]{ .949,  .949,  .949} rStar\cite{rSTaR} & Exploration & Human-Designed & 5 & ICL \& RL & \textcolor{green}{\ding{51}} & \encircle[fill=lightcoral, text=white]{T} & VC, GO, EB \\
         
        LLaMA-Berry\cite{zhang2024llama} & Exploration & Human-Designed & 2 & ICL \& RL & \textcolor{green}{\ding{51}} & \encircle[fill=lightcoral, text=white]{T} & VC, EB \\
    
        \rowcolor[rgb]{ .949,  .949,  .949} Huatuo-o1\cite{Huatuo-o1} & Synthetic Data & Human-Designed & 4 & ICL \& SFT & \textcolor{green}{\ding{51}} & \encircle[fill=lightcoral, text=white]{T} & VC \\

        Marco-o1\cite{Marco_o1} & Verification & Human-Designed & 1 & ICL \& SFT & \textcolor{green}{\ding{51}} & \encircle[fill=lightcoral, text=white]{T} & VC \\
        
         \rowcolor[rgb]{ .949,  .949,  .949}BoT\cite{yang2024buffer} & Exploration & In-Context Learning & Dynamic & ICL & \textcolor{red}{\ding{55}} & \encircle[fill=lightcoral, text=white]{T} & AD, IPR \\
         
        rStar-Math\cite{guan2025rstarmathsmallllmsmaster} & Exploration &  In-Context Learning & $1$ & ICL \& RL & \textcolor{green}{\ding{51}} & \encircle[fill=lightcoral, text=white]{T} & AD, IPR \\
        
         \rowcolor[rgb]{ .949,  .949,  .949}Mulberry\cite{yao2024mulberry} & Synthetic Data & In-Context Learning & 1 & ICL \& SFT & \textcolor{green}{\ding{51}}  & \encircle[fill=lightcoral, text=white]{T} & VC, EB \\
         
        LLaVA-CoT\cite{xu2024llava} & Synthetic Data/Exploration & Human-Designed & 4 & SFT & \textcolor{red}{\ding{55}} & \encircle[fill=harvestgold, text=white]{I} \space \encircle[fill=lightcoral, text=white]{T} & AD, IPR \\
        
         \rowcolor[rgb]{ .949,  .949,  .949}LLaMAV-o1\cite{thawakar2025llamav} & Verification/Exploration & Human-Designed & 4173 & Curriculum Learning & \textcolor{green}{\ding{51}} &  \encircle[fill=harvestgold, text=white]{I} \space \encircle[fill=lightcoral, text=white]{T} & AD, IPR \\
         
        AtomThink\cite{atomthink} & Synthetic Data/Exploration & In-Context Learning & $>$100 & SFT \& RL & \textcolor{green}{\ding{51}} & \encircle[fill=harvestgold, text=white]{I} \space \encircle[fill=lightcoral, text=white]{T} & AD, IPR, EB \\

         \rowcolor[rgb]{ .949,  .949,  .949}RedStar\cite{RedStar} & Distill & Human-Designed & 2 & SFT & \textcolor{green}{\ding{51}} & \encircle[fill=harvestgold, text=white]{I} \space \encircle[fill=lightcoral, text=white]{T} & AD, VC \\
        
        Auto-CoT \cite{zhangautomatic}  & Exploration & In-Context Learning & 2 & ICL & \textcolor{red}{\ding{55}} & \encircle[fill=lightcoral, text=white]{T} & AD, IPR, GO \\
        
         \rowcolor[rgb]{ .949,  .949,  .949}PoT\cite{chen2022program} & Verification & In-Context Learning & 1 & ICL & \textcolor{red}{\ding{55}} & \encircle[fill=lightcoral, text=white]{T} & AD, IPR, GO \\
        
        PAL\cite{gao2023pal} & Verification & In-Context Learning & 1 & ICL & \textcolor{red}{\ding{55}} & \encircle[fill=lightcoral, text=white]{T} & AD, IPR, GO \\
        
         \rowcolor[rgb]{ .949,  .949,  .949} Decomposed Prompt\cite{decomposed} & Exploration & Human-Designed & 3 & ICL & \textcolor{red}{\ding{55}} & \encircle[fill=lightcoral, text=white]{T} & AD, IPR \\
        
        Least-to-Most\cite{least2most} & Exploration & Human-Designed & 2 & ICL & \textcolor{red}{\ding{55}} & \encircle[fill=lightcoral, text=white]{T} & AD, IPR \\
        \bottomrule
   \end{tabular}
}
    \label{table:macro_action}
\end{table*}










% \begin{table*}[t!]
% \centering
% \renewcommand\arraystretch{1.2}
% \caption{Summary of Macro Action technology, classified into Action-guided Inference and Action-guide Data Synthesis.}
% % \bgroup
% % \def\arraystretch{1,2}
% \resizebox{0.98\linewidth}{!}{
% \begin{tabular}{llll}
% \toprule[1.2pt]
%   & \textbf{Reasoning LLMs}   & \textbf{Characteristic}           & \textbf{Benchmark}                          \\ \hline
  
% \fontsize{9}{6}\selectfont\multirow{5}{*}{\rotatebox{90}{{{\textbf{Inference}}}}} &  \cellcolor[rgb]{ .949,  .949,  .949}Self-Check \cite{DBLP:conf/iclr/MiaoTR24} &  \cellcolor[rgb]{ .949,  .949,  .949}Test-time Self-Improvement by Special Macro Action  &  \cellcolor[rgb]{ .949,  .949,  .949}GSM8K, MathQA, MATH \\

% & HiICL-MCTS \cite{HiICL-MCTS} & MCTS Search Enhanced by Macro Action Chain Prompt & GSM8K, MATH, StrategyQA, SVAMP  \\

% &  \cellcolor[rgb]{ .949,  .949,  .949}ReasonFlux \cite{ReasonFlux} &  \cellcolor[rgb]{ .949,  .949,  .949}Test-time Scaling by High-level Thought Iteration Process. &  \cellcolor[rgb]{ .949,  .949,  .949}MATH, AIME2024, AMC2023, OlympiadBench, Gaokao \\

% & rStar \cite{rSTaR} & MCTS Search Enhanced by Human-designed Macro Action  & GSM8K, GSM-Hard, SVAMP, StrategyQA    \\

% &  \cellcolor[rgb]{ .949,  .949,  .949}LLaMA-Berry \cite{DBLP:journals/corr/abs-2410-02884}  &  \cellcolor[rgb]{ .949,  .949,  .949}MCTS Search Enhanced by Human-designed Macro Action  &  \cellcolor[rgb]{ .949,  .949,  .949}GSM8K, MATH, GaoKao2023En, OlympiadBench, College Math, MMLU STEM  \\

% & Buffer of Thoughts \cite{yang2024buffer} & Inference by Retrieved High-level Thought.  & MGSM, BBH, Python Programming Puzzles, Shakespearean Sonnet Writing  \\

%  \hline

% \fontsize{9}{6}\selectfont\multirow{9}{*}{\rotatebox{90}{\textbf{Data Synthesis}}}

% &  \cellcolor[rgb]{ .949,  .949,  .949}rStar-Math \cite{guan2025rstarmathsmallllmsmaster} &  \cellcolor[rgb]{ .949,  .949,  .949}CoT Synthesis by Code Macro Action    &  \cellcolor[rgb]{ .949,  .949,  .949}MATH, AIME2024, AMC2023, OlympiadBench, College Math, GSM8K Gaokao   \\

% & Mulberry \cite{yao2024mulberry} & CoT Synthesis by Visual Macro Action  &  MathVista, MMStar, MMMU, ChartQA, DynaMath, HallBench, MM-Math MME$_{sum}$   \\

% &  \cellcolor[rgb]{ .949,  .949,  .949}LLaVA-CoT \cite{xu2024llava}  &  \cellcolor[rgb]{ .949,  .949,  .949}Human-designed Macro Action Tag Synthesis in CoT.  &  \cellcolor[rgb]{ .949,  .949,  .949}MMStar, MMBench, MMVet, MathVista, AI2D, Hallusion   \\

% & LLaMAV-o1 \cite{thawakar2025llamav} &  Structured Training Paradigm of Macro Action inserted CoT. & MMStar, MMBench, MMVet, MathVista, AI2D, Hallusion  \\

% &  \cellcolor[rgb]{ .949,  .949,  .949}AtomThink \cite{atomthink} &  \cellcolor[rgb]{ .949,  .949,  .949}Multimodal CoT Inserted Macro Action by ICL.  &  \cellcolor[rgb]{ .949,  .949,  .949}MathVista, MathVerse   \\

% & STILL-1  \cite{Scaling_of_Search_Learning} & Long-CoT Synthesis by Macro Action Assisted Tree Search & MATH-OAI, GSM-Hard, OlympiadBench, College Math   \\

% &  \cellcolor[rgb]{ .949,  .949,  .949}Satori  \cite{Satori} &  \cellcolor[rgb]{ .949,  .949,  .949}RL Explore by Macro Action Inserted Prompt  &  \cellcolor[rgb]{ .949,  .949,  .949}FOLIO, BGQA, CRUXEval, StrategyQA, TableBench, STEM    \\

% & RedStar \cite{RedStar}   & Distill Long-CoT by Macro Action Keyword Filter. &  MATH, AIME2024 OlympiadBench, College Math, High School Math    \\

% &  \cellcolor[rgb]{ .949,  .949,  .949}SUPERCORRECT \cite{yang2024supercorrect} &  \cellcolor[rgb]{ .949,  .949,  .949}Two-stage SFT Framework by High-level and Macro Action Thought Template. &  \cellcolor[rgb]{ .949,  .949,  .949}GSM8K, MATH, GaoKao \\

% \bottomrule[1.2pt]
% \end{tabular}
% }
% \label{table:macro_action}
% \end{table*}


%for \textit{System 2}

We classify the progress of macro action into two aspects: 
\begin{enumerate}[itemindent=0em]
\item \textbf{Test-time Scaling through Macro Action Operationalization:} Recent research identifies two key methodologies for improving reasoning performance during inference and test-time scaling. 
HiICL-MCTS \cite{HiICL-MCTS} employs a deliberate search through seed data to generate action-chain templates consisting of macro actions, thereby facilitating an action-chain-guided approach to test-time reasoning. ReasonFlux \cite{ReasonFlux} utilizes an iterative test-time scaling framework, harnessing external high-level thought templates to iteratively refine and update the current CoT.

%CoAct


\item \textbf{Macro Action-Enhanced Data Synthesis Paradigms:} A key application of macro actions in complex reasoning is in the synthesis of reasoning data. 
In data synthesis and training frameworks, macro action architectures enhance reasoning diversity and generalization. 
Recent research has shown that integrating or synthesizing a CoT process with macro actions within the reasoning sequence can significantly improve the data efficiency of the reasoning chain. 
For instance, LLaVA-CoT \cite{xu2024llava} enhances CoT data synthesis by externalizing intermediate reasoning steps across multiple modalities. 
AtomThink \cite{atomthink} generates the AMATH-SFT dataset using a structured g1 prompt \cite{g1}, achieving superior performance on long-horizon reasoning tasks compared to traditional CoT approaches. 
CoAct \cite{hou2024coact} introduces a dual-agent collaborative reasoning framework, where a global planning agent executes overarching macro-actions, while a local execution agent carries out specific sub-actions within those broader actions.

\end{enumerate}

Macro actions also play a crucial role in enhancing self improvement frameworks. 
rStar-Math \cite{guan2025rstarmathsmallllmsmaster} utilizes high-level deliberate search through Code-augmented CoT, generating diverse and reliable solutions while achieving proactive search capabilities. 
Satori \cite{Satori} integrates CoT with RL, incorporating ``$<$\textit{reflect}$>$''-style macro actions to diversify exploration and alleviate policy saturation in online RL environments. 
Huatuo-o1 \cite{Huatuo-o1} combines hierarchical planning with domain-specific knowledge bases to improve medical reasoning. 
Additionally, ReasonFlux \cite{ReasonFlux} dynamically reconfigures reasoning templates (\emph{e.g.}, breaking down calculus problems into symbolic and numeric phases) to align with the problem structure.





\subsubsection{Reinforcement Fine-Tuning}\label{rl_supervise}

Reinforcement Fine-Tuning (RFT) \cite{openaiRFT} is an innovative technique recently introduced by OpenAI, designed to enable developers and engineers to fine-tune existing models for specific domains or complex tasks. Unlike general SFT, RFT focuses on optimizing the model's reasoning process by using a reward mechanism to guide the model's evolution, thereby enhancing its reasoning capabilities and accuracy. 
The core of RFT lies in improving the model's performance in a specific domain with minimal high-quality training data \cite{chang2024survey}, an appropriate reward model \cite{trung2024reft}, and a stable optimization process in long-context \cite{blobRFT, codepmp, nextlong, quest}. 
A summary of RFT method is presented in Table \ref{table:rl_supervise}.







DeepSeek-R1 \cite{Deepseek-R1}, which employs a verifier reward-based strategy, has shown significant performance improvements compared to traditional methods like SoS \cite{SoS}. Key advantages include:
\begin{enumerate}[itemindent=0em]
    \item \textbf{Simplified Training Pipeline:} RL supervision streamlines data construction and training processes, eliminating the need for complex stepwise search mechanisms. 
    
    \item \textbf{Enhanced Scalability:} Online RL training facilitates efficient scaling on large datasets, particularly for complex reasoning tasks.
    
    \item \textbf{Emergent Properties:} DeepSeek-R1 \cite{Deepseek-R1} demonstrates unique emergent capabilities, such as Long-CoT reasoning, which are difficult to achieve through SFT alone.
\end{enumerate}


% \begin{table*}[tbp]
% \centering
% \renewcommand\arraystretch{1.3}
% \caption{Summary of  technology, including Reproduction work and Analysis work.}
% \resizebox{0.98\textwidth}{!}{
% \begin{tabular}{lllp{32em}l} % 修正列布局
% \toprule[1.2pt]
%     & \textbf{Reasoning LLMs} & \textbf{Characteristic} & \textbf{Benchmark} \\
% \hline
%        \fontsize{8}{6}\selectfont\multirow{11}{*}{\rotatebox{90}{\textbf{Reproduction}}}
%     &\cellcolor[rgb]{ .949,  .949,  .949}DeepSeek-R1 \cite{Deepseek-R1} & \cellcolor[rgb]{ .949,  .949,  .949}Large-scale RL, GRPO, Distillation & \cellcolor[rgb]{ .949,  .949,  .949}ArenaHard, LiveCodeBench, Codeforces, AIME 2024, MATH-500, etc. \\ 

%     %DROP, IF-Eval, GPQA, SimpleQA, FRAMES, AlpacaEval2.0, ArenaHard, LiveCodeBench, Codeforces, SWE Verified, Aider-Polyglot, AIME 2024, MMLU, MATH-500, CNMO 2024, CLUEWSC, C-Eval, C-SimpleQA
    
%     & ReFT \cite{trung2024reft} & RFT, PPO & GSM8k, SVAMP, MathQA\\
    
%     & \cellcolor[rgb]{ .949,  .949,  .949}T1 \cite{advancinglanguagemodelreasoning} & \cellcolor[rgb]{ .949,  .949,  .949}SFT, RLHF, RL scaling & \cellcolor[rgb]{ .949,  .949,  .949}MATH-500, AIME 2024, Omni-math-500 \\
    
%     & Satori \cite{Satori} & Autoregressive searching, Chain-of-Action-Thought &MATH-500, AIME 2024, CRUXEval, StrategyQA, TableBench, etc. \\

%     %GSM8K, MATH500, OlympiadBench, AMC, AIME 2024, FOLIO, BGQA, CRUXEval, StrategyQA, TableBench, STEM
    
%     & \cellcolor[rgb]{ .949,  .949,  .949}RLSP \cite{ye2025emergencethinkingllmsi} & \cellcolor[rgb]{ .949,  .949,  .949}RL via Self-play, SFT, PPO & \cellcolor[rgb]{ .949,  .949,  .949}MATH-500, AIME 2024 \\
    
%     & QCLASS \cite{QCLASS} & Exploration tree, PRM, Q-guided Generation Strategy & WebShop, SciWorld, ALFWorld \\
    
%     & \cellcolor[rgb]{ .949,  .949,  .949}OREAL \cite{OREAL} & \cellcolor[rgb]{ .949,  .949,  .949}Behavior Cloning and Reward Shaping, Outcome Reward-based RL & \cellcolor[rgb]{ .949,  .949,  .949}MATH-500, AIME 2024, AIME 2025-I, LiveMath, Olympiad \\
    
%     & DHP \cite{DHP} & Discrete Hierarchical Planning, Tree Trajectory Advantage Estimation & 25-room task \\
    
%     & \cellcolor[rgb]{ .949,  .949,  .949}DuoGuard \cite{DuoGuard} & \cellcolor[rgb]{ .949,  .949,  .949}Adversarial RL,  Multilingual Guardrail Modeling & \cellcolor[rgb]{ .949,  .949,  .949}XSTest, ToxicChat, OpenAI Moderation, Beavertails, RTP-LX, XSafety \\
    
%     & PRIME \cite{cui2025process} & Implicit PRM, RL & AIME 2024, AMC, Minerva Math, OlympiadBench, MATH-500 \\
    
%     & \cellcolor[rgb]{ .949,  .949,  .949}DeepScaleR \cite{deepscaler2025} & \cellcolor[rgb]{ .949,  .949,  .949}RL Scaling for Small Models, Iterative Lengthening for Length Scaling & \cellcolor[rgb]{ .949,  .949,  .949}AIME 2024, AMC, Minerva Math, OlympiadBench, MATH-500 \\
    
% \hline

%     \fontsize{8}{6}\selectfont\multirow{4}{*}{\rotatebox{90}{\textbf{Analysis}}} & DOES RLHF SCALE? \cite{RLHF_Scaling} & Analysis of RLHF Scaling Properties & MATH, GPQA, LiveCodeBench, GSM8K, MMLU \\
    
%     & \cellcolor[rgb]{ .949,  .949,  .949}Demystifying Long CoT Reasoning \cite{yeo2025demystifying} & \cellcolor[rgb]{ .949,  .949,  .949}Analysis of the Mechanics of Long CoT Reasoning using RL & \cellcolor[rgb]{ .949,  .949,  .949}MATH-500, AIME 2024, TheoremQA, MMLU-Pro \\
    
%     & Metastable Dynamics of CoT Reasoning \cite{metastabledynamicschainofthoughtreasoning} & 
% Analysis of Inference-time Compute and CoT Generation as a Metastable Markov Process & 
%     - \\
    
%     & \cellcolor[rgb]{ .949,  .949,  .949}Safety in DeepSeek-R1 \cite{parmar2025challenges} & \cellcolor[rgb]{ .949,  .949,  .949}Analysis of Limitations of RL for Reducing Harmful Outputs in DeepSeek-R1 & \cellcolor[rgb]{ .949,  .949,  .949}
%     - \\
% \bottomrule[1.2pt]
% \end{tabular}
% }
% \label{table:rl_supervise}
% \end{table*}

\begin{table*}[!t]
\centering
\caption{Summary of RFT method. \textbf{Tag}: \encircle[fill=harvestgold, text=white]{I} = \underline{I}mage, \encircle[fill=lightcoral, text=white]{T} = \underline{T}ext, \encircle[fill=DarkGreen, text=white]{V} = \underline{V}ideo.}
% We summarize the key features for each work: the \textbf{Foundation}(Foundation Model) and \textbf{Modality} of the feedback, the \textbf{Reward Type}, \textbf{Algorithm}, \textbf{Learning},  and \textbf{Incentivize Sample} size of the refinement process, and the \textbf{Application \& Benchmark} of the method. 
\resizebox{0.98\textwidth}{!}{
   \begin{tabular}% 
   {llclllll}
    \toprule
        \multirow{2}{*}{\textbf{Methods}} & \multicolumn{2}{c}{\textbf{Model Attribute}} & \multicolumn{4}{c}{\textbf{Incentivize Attibute}} & \multirow{2}{*}{\textbf{Application \& Benchmark}} \\ \cline{2-3} \cline{4-7}
         & \textbf{Foundational LLMs} & \textbf{Feedback Modality} & \textbf{Reward Type} & \textbf{Algorithm} & \textbf{Learning} & \textbf{Incentivize Sample} & \\ \midrule
        \multicolumn{8}{c}{\textbf{Reason RFT Project}} \\ \midrule
        % RLHF
         \rowcolor[rgb]{ .949,  .949,  .949}DeepSeek-R1-Zero\cite{Deepseek-R1} & DeepSeek-V3 & \encircle[fill=lightcoral, text=white]{T} & Rule-Outome-Reward & GPRO & RL & 800K  & Multiple Tasks \\
        
        DeepSeek-R1\cite{Deepseek-R1} & DeepSeek-V3 & \encircle[fill=lightcoral, text=white]{T} & Rule-Outcome-Reward & GPRO & RL \& SFT & 800K  & Multiple Tasks \\
        
         \rowcolor[rgb]{ .949,  .949,  .949}Kimi v1.5\cite{team2025kimi} & -- & \encircle[fill=harvestgold, text=white]{I} \space\encircle[fill=lightcoral, text=white]{T} & Rule-Outcome-Reward & $\text{PPO}^{*}$ & RL \& SFT & -- & Multiple Tasks \\
         
        % Fine-grained RLHF
        ReFT\cite{trung2024reft}  & Galactica, CodeLLama & \encircle[fill=lightcoral, text=white]{T} & Rule-Outcome-Reward & $\text{PPO}^{*}$ & RL \& SFT & 3k/7k/8k/15k & GSM8k/SVAMP/MathQA \\
        
         \rowcolor[rgb]{ .949,  .949,  .949}RFTT\cite{RFTT} & LLaMA-3-3/8B-Instruct,Qwen-2.5-7B-Instruct & \encircle[fill=lightcoral, text=white]{T} & Rule-Outcome-Reward &  Reinforce++ & RL \& SFT & 1.2K & Multiple Math Task \\
         
        % HH-RLHF
        Satori\cite{Satori} & Qwen-2.5-Math-7B & \encircle[fill=lightcoral, text=white]{T} & Rule-Outcome-Reward & PPO & RL \& SFT & 66K & Multiple Math Task  \\
        
         \rowcolor[rgb]{ .949,  .949,  .949}QCLASS\cite{QCLASS} & Llama-2-7B-Chat & \encircle[fill=lightcoral, text=white]{T} & Process-Reward & QNet & RL \& SFT & 1.9K/1.5K/3.3K & WebShop, ALFWorld, SciWorld \\
        
        % ILF
        PRIME\cite{cui2025process} & Qwen2.5-Math-7B & \encircle[fill=lightcoral, text=white]{T} & Rule-Process-Outcome-Reward & PPO & RL \& SFT  & 150K & Math, Code Tasks  \\
        
         \rowcolor[rgb]{ .949,  .949,  .949}DeepScaleR\cite{deepscaler2025} & DeepSeek-R1-Distill-Qwen-1.5B & \encircle[fill=lightcoral, text=white]{T} & Rule-Outcome-Reward & Iteratively GPRO & RL & 40K & Multiple Math Task \\
        
        PURE\cite{cheng2025pure} & Qwen2.5-Math-7B &  \encircle[fill=lightcoral, text=white]{T} & Rule-Process-Outcome-Reward & PPO+RLOO & RL & 8K & Multiple Math Task \\
        
         \rowcolor[rgb]{ .949,  .949,  .949}SimpleRL\cite{simplerl_reason_blob} & Qwen2.5-Math-7B & \encircle[fill=lightcoral, text=white]{T} & Rule-Outcome-Reward & PPO & RL & 8K & Multiple Math Task \\
        
        Open-R1\cite{openr1} & Qwen2.5-1.5B-Instruct & \encircle[fill=lightcoral, text=white]{T} & Rule-Outcome-Reward & GPRO & RL \& SFT & 8K & Multiple Math, Code Task \\
        
         \rowcolor[rgb]{ .949,  .949,  .949}TinyZero\cite{tinyzero} & Qwen2.5-0.5B/3B & \encircle[fill=lightcoral, text=white]{T} & Rule-Outcome-Reward & GPRO & RL & -- & CountDown Task \\
        
        Ota-Zero\cite{liu2025oatzero} & Qwen-2.5-Series, DeepSeek-Series, Rho, Llama-3.x & \encircle[fill=lightcoral, text=white]{T} & Rule-Outcome-Reward & GRPO & RL & 0.5K & CountDown Task \\
        
         \rowcolor[rgb]{ .949,  .949,  .949}Ota\cite{liu2025oat} & RHO-1b/Qwen2.5-3B & \encircle[fill=lightcoral, text=white]{T} & Rule-Outcome-Reward & GPRO/PPO & RL & 7.5K & GSM8K \\
        
        LIMR\cite{limr} & Qwen-Math-7B & \encircle[fill=lightcoral, text=white]{T} & Rule-Outcome-Reward & PPO & RL & 1.3K & Multiple Math Task \\
        
         \rowcolor[rgb]{ .949,  .949,  .949}Critic-RL\cite{critic_rl} & Qwen2.5-Coder-32B & \encircle[fill=lightcoral, text=white]{T} & Rule-Outcome-Reward & $\text{GPRO}^{*}$ & RL \& SFT & 18.8K & Multiple Code Task \\
        
        Logic-R1\cite{logicrl} & Qwen2.5-7B-Instruct-1M & \encircle[fill=lightcoral, text=white]{T} & Rule-Outcome-Reward & $\text{REINFORCE++}^{*}$ & RL & 5K & Multiple Math, Logic Task \\
        
         \rowcolor[rgb]{ .949,  .949,  .949}Online-DPO-R1\cite{online_dpo_r1} & Qwen2.5-MATH-7B & \encircle[fill=lightcoral, text=white]{T} & Rule-Outcome-Reward  & DPO  & RL\& SFT & 207.5K & Multiple Math Task \\
        
        OpenReason-Zero\cite{OpenReasonerZero2025} & Qwen-2.5-7B/32B & \encircle[fill=lightcoral, text=white]{T} & Rule-Outcome-Reward   & PPO & RL & 57K & Multiple Math Task, GPQA, MMLU \\
        
         \rowcolor[rgb]{ .949,  .949,  .949}RLHF-V\cite{yu2024rlhfv} & OmniLMM-12B & \encircle[fill=harvestgold, text=white]{I} \space\encircle[fill=lightcoral, text=white]{T} & Process-Reward &  DDPO & RL & 1.4K & Multiple Tasks \\
        
        RLAIF\cite{lee2023rlaif} & PaLM 2 Extra-Small & \encircle[fill=lightcoral, text=white]{T} & Rule-Outome-Reward & RLAIF & RL & -- & Summary and Conversation Generation \\
        
         \rowcolor[rgb]{ .949,  .949,  .949}MM-RLHF\cite{MM_RLHF} & LLaVA-onevision-7B & \encircle[fill=harvestgold, text=white]{I} \space\encircle[fill=lightcoral, text=white]{T} \space\encircle[fill=DarkGreen, text=white]{V} & Process-Reward & MM-DPO & RL & 120K & MM-RLHF-RewardBench/SafetyBench\\
        
         Align-DS-V\cite{Align_DS_V}  & LLaVA-v1.5-7B,Qwen2-VL &\encircle[fill=harvestgold, text=white]{I} \space\encircle[fill=lightcoral, text=white]{T} \space\encircle[fill=DarkGreen, text=white]{V}
        & Process-Reward & PPO, DPO 
        & RL \& SFT & 200K & Align-Anything, Eval-Anything \\
        
         \rowcolor[rgb]{ .949,  .949,  .949}R1V\cite{chen2025r1v}  & Qwen2-VL,Qwen2.5-VL
        & \encircle[fill=harvestgold, text=white]{I} \space\encircle[fill=lightcoral, text=white]{T} 
        &  Rule-Outome-Reward & GRPO & RL & 70K/70K/8K & Multiple Tasks \\
        
        % HH-RLHF
        VLM-R1\cite{vlmr1} &  Qwen2.5-VL & \encircle[fill=harvestgold, text=white]{I} \space\encircle[fill=lightcoral, text=white]{T}  
        & Rule-Outome-Reward  & GRPO & RL  & 120K & Multiple Tasks \\
        
         \rowcolor[rgb]{ .949,  .949,  .949}LMM-R1\cite{lmmr1} & Qwen2.5-VL&\encircle[fill=harvestgold, text=white]{I} \space\encircle[fill=lightcoral, text=white]{T} 
        &  Rule-Outome-Reward  &  PPO/RLOO & RL  & 8K & Multiple Tasks \\
        
        Open-R1-Video\cite{open-r1-video} & Qwen2-VL-7B  & \encircle[fill=harvestgold, text=white]{I} \space\encircle[fill=lightcoral, text=white]{T} \space\encircle[fill=DarkGreen, text=white]{V}
        & Rule-Outome-Reward & GRPO  & RL & 4K & Multiple Tasks \\
        
         \rowcolor[rgb]{ .949,  .949,  .949}Easy-R1 \cite{zheng2025easyr1} & Qwen2.5-VL& \encircle[fill=harvestgold, text=white]{I} \space\encircle[fill=lightcoral, text=white]{T}
        &  Rule-Outome-Reward  & GRPO & RL & 3K & Multiple Tasks \\
        %%%%%%
        \midrule
        \multicolumn{8}{c}{\textbf{Analysis RFT Project}} \\ 
        \midrule
        %%%%%%
        % Self-Verification
        Demystify-LongCoT\cite{yeo2025demystifying} & Llama-3.1-8B, Qwen2.5
-7B-Math & \encircle[fill=lightcoral, text=white]{T} & Rule-Outcome-Reward  & PPO/Reinforce++ &  RL \& SFT  & 7.5K & Multiple Math, MMLU \\ 

         \rowcolor[rgb]{ .949,  .949,  .949}RLHF-Scale\cite{RLHF_Scaling} & GLM4-9B & \encircle[fill=lightcoral, text=white]{T} & Process-Reward  & PPO & RL  & 11K & Multiple Tasks  \\ 
        
        MD-CoT\cite{metastabledynamicschainofthoughtreasoning}  & -- & -- & Process-Reward  & PPO & RL  & -- & --  \\ 
        
        %%%%%%
        \bottomrule
   \end{tabular}
}
    \label{table:rl_supervise}
\end{table*}



Despite its strengths, RFT faces the following challenges:

\begin{enumerate}[itemindent=0em]

\item \textbf{Unclear Mechanism behind Reasoning:} The underlying mechanisms driving the reasoning improvements in DeepSeek-R1 remain poorly understood. 
For example, while DeepSeek-R1 exhibits emergent properties (\emph{e.g.}, ``Emergent Length Increasing'', ``Aha moments''), studies such as \cite{ThereMaybeNot} suggest that capabilities like Long-CoT might already exist in the base model, rather than solely emerging from RL training. 
Furthermore, performance gains observed in smaller models (\emph{e.g.}, Qwen-Math-2B/7B \cite{qwen2}) occur without noticeable ``Aha moments'', complicating causal interpretations.

\item \textbf{Reward Model Saturation:} Many existing RL algorithms face reward model saturation, typically manifested as exploration collapse after around 100 training steps. 
Although DeepSeek-R1 alleviates this issue through specialized reward formatting, methods like ReFT \cite{trung2024reft} and Satori \cite{Satori} propose alternating sampling and SFT distillation to combat reward hacking and exploration collapse.

\item \textbf{Unstable Long-CoT Generation:} Long reasoning chains generated by RFT are prone to instability, including context overflow, failure to return final answers, and sensitivity to reward shaping \cite{Tecent_2_plus_3}. 
For instance, methods like \cite{yeo2025demystifying} inadvertently introduce cosine reward functions, which degrade performance with increased iterations. 
O1-Prune \cite{o1_pruner} uses post-hoc length pruning techniques \cite{team2025kimi} (via RL/SFT) to stabilize outputs.

\end{enumerate}

% 技术介绍 优势，不足，未来发展 有代表性的方法

Future directions for RFT may include several exciting and innovative advancements, such as:

\begin{enumerate}[itemindent=0em]

\item \textbf{Efficient and Stable RL Frameworks:} There is a need to develop more robust RL algorithms that prevent reward saturation and exploration collapse. \cite{yeo2025demystifying} reveals that REINFORCE++ \cite{reinforce_plusplus} underperforms when combined with KL divergence regularization, suggesting the need for alternative methods. 
Future work should revisit classic RL algorithms in the context of modern LLMs training to optimize both stability and efficiency.

\item \textbf{Scaling RFT:} Current RL-Supervise models rely on curated, verifiable prompts selected from large-scale datasets. 
Future research should focus on synthesizing high-quality, diverse prompts to improve generalization. \cite{RLHF_Scaling} shows that merely scaling policy/reward models or increasing sample sizes results in diminishing returns, while expanding the scope of PRM and R1 training data holds greater promise. 
Hybrid approaches, such as combining RL with SFT or curriculum learning, should be explored to enhance scalability.

\item \textbf{Controlling Long-CoT Stability:} Adaptive reward shaping mechanisms are needed to balance reasoning length, coherence, and answer correctness. 
Techniques such as O1-Prune \cite{o1_pruner} demonstrate the value of post-hoc length regularization, but dynamic in-training controls are necessary. 
Hierarchical RL frameworks should be investigated to decompose long reasoning chains into manageable sub-tasks, reducing instability.

\item \textbf{Theoretical and Empirical Analysis:} It is essential to clarify the relationship between RL training and the capabilities of the base model. 
For instance, it should be determined whether emergent properties (\emph{e.g.}, Long-CoT) arise from RL optimization or are latent traits of the base model. Systematic studies on reward design principles (\emph{e.g.}, sparse vs. dense rewards, multi-objective balancing) should be conducted to avoid unintended behaviors such as reward hacking.

\end{enumerate}

\noindent\textbf{Summary:} RFT presents a promising direction for advancing LLMs reasoning, as evidenced by DeepSeek-R1 \cite{Deepseek-R1}. 
However, challenges such as reward saturation, unstable long reasoning chains, and unclear emergent mechanisms require urgent attention. 
Future efforts should prioritize algorithmic innovation, scalable prompt synthesis, and theoretical grounding to fully unlock the potential of RL-driven reasoning LLMs.





\subsection{Evolutionary of Reasoning LLMs}\label{evolutionary}

% In this section, we summarize the various developmental stages and the evolutionary timeline of LLMs-driven SRS.

The evolution of reasoning LLMs has progressed by several distinct stages, with various strategies developed to overcome the limitations of direct autoregressive inference and build more advanced slow-thinking reasoning architectures.

% The evolutionary timeline of Reasoning LLMs is presented in Figure \ref{fig:timeline}.


% \begin{table*}[tbhp]
% \scriptsize
% \centering
% \caption{Statistics of benchmarks for reasoning LLMs.}
% \resizebox{0.91\linewidth}{!}{
% \begin{tabular}{llllll}
% \toprule[1.2pt]
% \textbf{Domain}                         & \textbf{Benchmark}                                                            & \textbf{Venue}      & \textbf{Language}        & \textbf{Size}   & \textbf{Level}              \\
% \hline
% \multirow{4}{*}{Math}         & \cellcolor[rgb]{ .949,  .949,  .949} AIME 2024 \cite{AIME2024}                           & \cellcolor[rgb]{ .949,  .949,  .949}-          & \cellcolor[rgb]{ .949,  .949,  .949}English         &\cellcolor[rgb]{ .949,  .949,  .949} 30     & \cellcolor[rgb]{ .949,  .949,  .949}Competition        \\

%                                &  MATH-500 \cite{lightmanlet}                         &  ICLR 2024    & English         &  500    & Competition        \\
                               
%                                & \cellcolor[rgb]{ .949,  .949,  .949}AMC 2023 \cite{AMC2023}                             &\cellcolor[rgb]{ .949,  .949,  .949} -          & \cellcolor[rgb]{ .949,  .949,  .949}English         & \cellcolor[rgb]{ .949,  .949,  .949}30     & \cellcolor[rgb]{ .949,  .949,  .949}Competition        \\
                               
%                                &  Olympiad Bench   \cite{he2024olympiadbench}         &  ACL 2024     & English/Chinese &  8,476  & Competition        \\
% \hline
% \multirow{3}{*}{Code}          & \cellcolor[rgb]{ .949,  .949,  .949}Codeforces                                                           & \cellcolor[rgb]{ .949,  .949,  .949}-          & \cellcolor[rgb]{ .949,  .949,  .949}English         & \cellcolor[rgb]{ .949,  .949,  .949}-      &\cellcolor[rgb]{ .949,  .949,  .949} Expert             \\

%                                &  SWE-bench   \cite{jimenez2024swebench}              &  ICLR 2024    &  English         &  2,294  &  Expert             \\
                               
%                                & \cellcolor[rgb]{ .949,  .949,  .949}LiveCodeBench   \cite{jain2024livecodebench}        &\cellcolor[rgb]{ .949,  .949,  .949} ArXiv 2024   & \cellcolor[rgb]{ .949,  .949,  .949}English         & \cellcolor[rgb]{ .949,  .949,  .949}-      & \cellcolor[rgb]{ .949,  .949,  .949}Expert             \\
                               
% \hline
% \multirow{3}{*}{Science}       &  GPQA Diamond \cite{rein2024gpqa}                    &  COLM 2024    &  English         & 448    &  University         \\

%                                & \cellcolor[rgb]{ .949,  .949,  .949}MR-Ben \cite{ZengLWLCDYXQZSL24}                        & \cellcolor[rgb]{ .949,  .949,  .949}NeurIPS 2024 & \cellcolor[rgb]{ .949,  .949,  .949}English         & \cellcolor[rgb]{ .949,  .949,  .949}5,975 & \cellcolor[rgb]{ .949,  .949,  .949} Hybrid             \\
%                                & \cellcolor[rgb]{ .949,  .949,  .949}MMLU-Pro \cite{wang2024mmlu}                        & \cellcolor[rgb]{ .949,  .949,  .949}NeurIPS 2024 & \cellcolor[rgb]{ .949,  .949,  .949}English         & \cellcolor[rgb]{ .949,  .949,  .949}12,032 & \cellcolor[rgb]{ .949,  .949,  .949}Hybrid             \\
                               
% \hline

% \multirow{4}{*}{Agent}         &  WebShop \cite{yao2022webshop}                       &  NeurIPS 2022 &  English         &  1,600  & Hybrid             \\

%                                & \cellcolor[rgb]{ .949,  .949,  .949} WebArena   \cite{zhou2023webarena}                  & \cellcolor[rgb]{ .949,  .949,  .949}ICLR 2024    & \cellcolor[rgb]{ .949,  .949,  .949}English         & \cellcolor[rgb]{ .949,  .949,  .949}812    & \cellcolor[rgb]{ .949,  .949,  .949}Hybrid             \\
                               
%                                &  SciWorld   \cite{scienceworld2022}                  &  EMNLP 2022   & English         & 7,200  &  Hybrid             \\
                               
%                                & \cellcolor[rgb]{ .949,  .949,  .949}TextCraft   \cite{prasad2024adapt}           & \cellcolor[rgb]{ .949,  .949,  .949}NAACL 2024   & \cellcolor[rgb]{ .949,  .949,  .949}English         & \cellcolor[rgb]{ .949,  .949,  .949}200    & \cellcolor[rgb]{ .949,  .949,  .949}Hybrid             \\
                               
% \hline

% \multirow{3}{*}{Medicine}      &  JAMA Clinical Challenge \cite{chen2024benchmarking} &  NAACL 2025   & English         &  1,524  &  Expert             \\

%                                & \cellcolor[rgb]{ .949,  .949,  .949}Medbullets   \cite{chen2024benchmarking}            & \cellcolor[rgb]{ .949,  .949,  .949}NAACL 2025   & \cellcolor[rgb]{ .949,  .949,  .949}English         & \cellcolor[rgb]{ .949,  .949,  .949}308    & \cellcolor[rgb]{ .949,  .949,  .949}Expert             \\
                               
%                                &  MedQA \cite{jin2021disease}                         &  ArXiv 2020   & English/Chinese &  61,097 &  Expert             \\
                               
% \hline
% \multirow{5}{*}{Multimodality} & \cellcolor[rgb]{ .949,  .949,  .949}MMMU \cite{yue2024mmmu}                             & \cellcolor[rgb]{ .949,  .949,  .949}CVPR 2024    & \cellcolor[rgb]{ .949,  .949,  .949}English         & \cellcolor[rgb]{ .949,  .949,  .949}11,500 & \cellcolor[rgb]{ .949,  .949,  .949}Hybrid             \\

%                                &  MathVista   \cite{lu2024mathvista}                  &  ICLR 2024    &  English         &  6,141  &  Middle School      \\
                               
%                                & \cellcolor[rgb]{ .949,  .949,  .949}MathVision \cite{MathVision}                        & \cellcolor[rgb]{ .949,  .949,  .949}NeurIPS 2024 & \cellcolor[rgb]{ .949,  .949,  .949}English         & \cellcolor[rgb]{ .949,  .949,  .949}3,040  & \cellcolor[rgb]{ .949,  .949,  .949}Middle/High School \\
                               
%                                &  CMMaTH \cite{li2024cmmath}                          &  COLING 2025  & English/Chinese &  23,856 & Middle/High School \\
                               
%                                & \cellcolor[rgb]{ .949,  .949,  .949}PGPS9K \cite{Zhang2023PGPS}                         & \cellcolor[rgb]{ .949,  .949,  .949}IJCAI 2023   & \cellcolor[rgb]{ .949,  .949,  .949}English         & \cellcolor[rgb]{ .949,  .949,  .949}9,023  & \cellcolor[rgb]{ .949,  .949,  .949}Middle School  \\ 
% \bottomrule[1.2pt]
% \end{tabular}
% }
% \label{table:benchmark_categories}
% \end{table*}



\begin{table*}[tbhp]
\scriptsize
\centering
\caption{Statistics of benchmarks for reasoning LLMs.}
\resizebox{0.95\linewidth}{!}{
\begin{tabular}{lllllll}
\toprule[1.2pt]
\textbf{Domain}                  & \textbf{Benchmark}                                                            & \textbf{Question Type} & \textbf{Venue}      & \textbf{Language}        & \textbf{Size}   & \textbf{Level}              \\
\hline
\multirow{4}{*}{Math}         

& \cellcolor[rgb]{ .949,  .949,  .949} AIME 2024 \cite{AIME2024}                        &  \cellcolor[rgb]{ .949,  .949,  .949} Open-End & \cellcolor[rgb]{ .949,  .949,  .949}-          & \cellcolor[rgb]{ .949,  .949,  .949}English         &\cellcolor[rgb]{ .949,  .949,  .949} 30     & \cellcolor[rgb]{ .949,  .949,  .949}Competition        \\

&  MATH-500 \cite{lightmanlet}  & Open-End &  ICLR 2024 & English & 500  & Competition        \\

& \cellcolor[rgb]{ .949,  .949,  .949} AMC 2023 \cite{AMC2023}                          & \cellcolor[rgb]{ .949,  .949,  .949} Open-End & \cellcolor[rgb]{ .949,  .949,  .949}--   & \cellcolor[rgb]{ .949,  .949,  .949}English         & \cellcolor[rgb]{ .949,  .949,  .949}30     & \cellcolor[rgb]{ .949,  .949,  .949}Competition        \\

&  Olympiad Bench   \cite{he2024olympiadbench}    & Open-End  &  ACL 2024     & English/Chinese &  8,476  & Competition        \\
\hline

\multirow{3}{*}{Code}          

& \cellcolor[rgb]{ .949,  .949,  .949}Codeforces                                         & \cellcolor[rgb]{ .949,  .949,  .949} Open-End   & \cellcolor[rgb]{ .949,  .949,  .949}-          & \cellcolor[rgb]{ .949,  .949,  .949}English         & \cellcolor[rgb]{ .949,  .949,  .949}-      &\cellcolor[rgb]{ .949,  .949,  .949} Expert             \\

&  SWE-bench \cite{jimenez2024swebench} & Open-End  &  ICLR 2024    &  English         &  2,294  &  Expert             \\

& \cellcolor[rgb]{ .949,  .949,  .949}LiveCodeBench   \cite{jain2024livecodebench}     & \cellcolor[rgb]{ .949,  .949,  .949} Open-End   &\cellcolor[rgb]{ .949,  .949,  .949} ArXiv 2024   & \cellcolor[rgb]{ .949,  .949,  .949}English         & \cellcolor[rgb]{ .949,  .949,  .949}-      & \cellcolor[rgb]{ .949,  .949,  .949}Expert             \\

\hline
\multirow{3}{*}{Science} 
&  GPQA Diamond \cite{rein2024gpqa}   & Choice &  COLM 2024    &  English         & 448    &  University  \\

& \cellcolor[rgb]{ .949,  .949,  .949} MR-Ben \cite{ZengLWLCDYXQZSL24}        & \cellcolor[rgb]{ .949,  .949,  .949} Hybrid & \cellcolor[rgb]{ .949,  .949,  .949}NeurIPS 2024 & \cellcolor[rgb]{ .949,  .949,  .949}English         & \cellcolor[rgb]{ .949,  .949,  .949}5,975 & \cellcolor[rgb]{ .949,  .949,  .949} Hybrid             \\

 & \cellcolor[rgb]{ .949,  .949,  .949}MMLU-Pro \cite{wang2024mmlu} & \cellcolor[rgb]{ .949,  .949,  .949}Choice  & \cellcolor[rgb]{ .949,  .949,  .949}NeurIPS 2024 & \cellcolor[rgb]{ .949,  .949,  .949}English         & \cellcolor[rgb]{ .949,  .949,  .949}12,032 & \cellcolor[rgb]{ .949,  .949,  .949}Hybrid             \\

\hline

\multirow{4}{*}{Agent}         
&  WebShop \cite{yao2022webshop}  & Open-End                   
&  NeurIPS 2022 &  English  &  1,600  & Hybrid \\

& \cellcolor[rgb]{ .949,  .949,  .949} WebArena  \cite{zhou2023webarena} & \cellcolor[rgb]{ .949,  .949,  .949}Open-End  & \cellcolor[rgb]{ .949,  .949,  .949}ICLR 2024    & \cellcolor[rgb]{ .949,  .949,  .949}English         & \cellcolor[rgb]{ .949,  .949,  .949}812    & \cellcolor[rgb]{ .949,  .949,  .949}Hybrid             \\

&  SciWorld   \cite{scienceworld2022}  & Open-End &  EMNLP 2022   & English   & 7,200  &  Hybrid  \\

& \cellcolor[rgb]{ .949,  .949,  .949}TextCraft   \cite{prasad2024adapt} & \cellcolor[rgb]{ .949,  .949,  .949}Open-End  & \cellcolor[rgb]{ .949,  .949,  .949}NAACL 2024   & \cellcolor[rgb]{ .949,  .949,  .949}English         & \cellcolor[rgb]{ .949,  .949,  .949}200    & \cellcolor[rgb]{ .949,  .949,  .949}Hybrid             \\

\hline
\multirow{3}{*}{Medicine}      
&  JAMA Clinical \cite{chen2024benchmarking} & Choice &  NAACL 2025   & English  &  1,524  &  Expert      \\

& \cellcolor[rgb]{ .949,  .949,  .949}Medbullets \cite{chen2024benchmarking}  & \cellcolor[rgb]{ .949,  .949,  .949}Choice & \cellcolor[rgb]{ .949,  .949,  .949}NAACL 2025   & \cellcolor[rgb]{ .949,  .949,  .949}English      & \cellcolor[rgb]{ .949,  .949,  .949}308    & \cellcolor[rgb]{ .949,  .949,  .949}Expert             \\

&  MedQA \cite{jin2021disease}  & Choice  &  ArXiv 2020   & English/Chinese &  61,097 &  Expert             \\

\hline
\multirow{5}{*}{Multimodality} 
& \cellcolor[rgb]{ .949,  .949,  .949}MMMU \cite{yue2024mmmu} & \cellcolor[rgb]{ .949,  .949,  .949}Hybrid & \cellcolor[rgb]{ .949,  .949,  .949}CVPR 2024    & \cellcolor[rgb]{ .949,  .949,  .949}English         & \cellcolor[rgb]{ .949,  .949,  .949}11,500 & \cellcolor[rgb]{ .949,  .949,  .949}Hybrid             \\

&  MathVista   \cite{lu2024mathvista}  & Hybrid &  ICLR 2024    &  English         &  6,141  &  Middle School      \\

& \cellcolor[rgb]{ .949,  .949,  .949}MathVision \cite{MathVision}   & \cellcolor[rgb]{ .949,  .949,  .949}Hybrid & \cellcolor[rgb]{ .949,  .949,  .949}NeurIPS 2024 & \cellcolor[rgb]{ .949,  .949,  .949}English         & \cellcolor[rgb]{ .949,  .949,  .949}3,040  & \cellcolor[rgb]{ .949,  .949,  .949}Middle/High School \\

&  CMMaTH \cite{li2024cmmath}  & Hybrid &  COLING 2025  & English/Chinese &  23,856 & Middle/High School \\

& \cellcolor[rgb]{ .949,  .949,  .949}PGPS9K \cite{Zhang2023PGPS}  & \cellcolor[rgb]{ .949,  .949,  .949}Hybrid & \cellcolor[rgb]{ .949,  .949,  .949}IJCAI 2023   & \cellcolor[rgb]{ .949,  .949,  .949}English         & \cellcolor[rgb]{ .949,  .949,  .949}9,023  & \cellcolor[rgb]{ .949,  .949,  .949}Middle School  \\ 
\bottomrule[1.2pt]
\end{tabular}
}
\label{table:benchmark_categories}
\end{table*}


In the early stages, reasoning LLMs primarily focused on enhancing pre-trained LLMs with external reasoning algorithms, without altering the underlying model parameters. Approaches such as Tree of Thoughts \cite{Tree_of_Thought} and Reasoning via Planning \cite{hao2023reasoning} utilized LLMs-driven Breadth-First Search, Depth-First Search, and MCTS \cite{Renda_Report_Tree_Search, browne2012survey, DBLP:journals/corr/abs-2410-16033, DBLP:journals/corr/abs-2412-09078} to simulate human-like reasoning processes. 
These methods represented reasoning as tree or graph traversals, where intermediate reasoning states were depicted as nodes, and various reasoning strategies produced distinct reasoning paths. The final decision was made through additional voting mechanisms \cite{wangself} or Monte Carlo-based value estimation to identify the optimal path.


However, these externalized slow-reasoning approaches introduced several challenges:
\begin{enumerate}[itemindent=0em]

\item \textbf{Limited Exploration Space:} The search-based methods required predefined constraints on the breadth, depth, and granularity of the search space, which often restricted the LLM's exploration to a narrow reasoning space. Furthermore, the reasoning strategies across different child nodes of the same parent node frequently lacked sufficient diversity, further limiting exploration.

\item \textbf{Limited Experience Sharing:} Exploration experiences and reasoning information across different paths could only be assessed based on reward models or self-consistency among outcomes. Additionally, search-based methods significantly increased computational overhead, relying on reward models such as PRM/ORM for tree pruning or speculative decoding techniques to accelerate inference.

\end{enumerate}
To overcome these limitations, subsequent models such as rSTaR \cite{rSTaR}, LLaMAV-o1 \cite{thawakar2025llamav}, HiICL-MCTS \cite{HiICL-MCTS}, Mulberry \cite{yao2024mulberry}, g1 \cite{g1}, and Thinking-Claude \cite{Thinking-Claude} introduced richer action spaces. 
These enhanced action spaces offered high-level planning cues, broadening the model's exploration scope and enabling more comprehensive structured search processes. 
However, this approach necessitated careful design of the action spaces to ensure their effectiveness.




With the introduction of models like o1 \cite{openai_o1} and QwQ \cite{qwq-32b-preview}, external reasoning paradigms were internalized within the LLM's context. 
These models initially performed exploratory macro-planning to generate an initial reasoning path, followed by contextual exploration of alternative paths. 
Through mechanisms like ``Rethink'' and ``Verification'', these models produced extended reasoning chains. 
To replicate this internalized capability, STILL-1 \cite{Renda_Report_Tree_Search} linearized tree search outputs into long reasoning chains with attributes such as ``Rethink'', ``Wait'', and ``Explore New Path''. 
Similarly, STILL-2 \cite{Slow_Thinking_with_LLMs_2} and sky-T1 \cite{sky_t1_2025} synthesized long reasoning chains using distillation techniques. 
However, the linearized reasoning chains derived from search-based methods struggled to match the quality of those produced by distillation approaches.





Recent advancements, including DeepSeek-R1 \cite{Deepseek-R1} and Kimi-k1.5 \cite{team2025kimi}, have demonstrated the potential of RL to enhance models like DeepSeek-V3 \cite{liu2024deepseek}, resulting in the emergence of complex behaviors such as long reasoning chains, reflective reasoning, and advanced planning capabilities. 
Remarkably, these sophisticated behaviors were achieved through simple RL scaling. 
SimpleRL \cite{simplerl_reason_blob} sought to replicate these capabilities using a streamlined pipeline and minimal codebase, while R1V \cite{chen2025r1v} explored the development of multimodal reasoning models based on multimodal foundation architectures.


\noindent\textbf{Summary:} The evolution of reasoning LLMs has shifted from externally augmented reasoning to internally embedded reasoning. 
Recent developments emphasize the potential of RL-based scaling to unlock advanced capabilities.



\begin{table}[t]
    \centering
    \caption{The performance of different pre-trained models on ImageNet and infrared semantic segmentation datasets. The \textit{Scratch} means the performance of randomly initialized models. The \textit{PT Epochs} denotes the pre-training epochs while the \textit{IN1K FT epochs} represents the fine-tuning epochs on ImageNet \citep{imagenet}. $^\dag$ denotes models reproduced using official codes. $^\star$ refers to the effective epochs used in \citet{iBOT}. The top two results are marked in \textbf{bold} and \underline{underlined} format. Supervised and CL methods, MIM methods, and UNIP models are colored in \colorbox{orange!15}{\rule[-0.2ex]{0pt}{1.5ex}orange}, \colorbox{gray!15}{\rule[-0.2ex]{0pt}{1.5ex}gray}, and \colorbox{cyan!15}{\rule[-0.2ex]{0pt}{1.5ex}cyan}, respectively.}
    \label{tab:benchmark}
    \centering
    \scriptsize
    \setlength{\tabcolsep}{1.0mm}{
    \scalebox{1.0}{
    \begin{tabular}{l c c c c  c c c c c c c c}
        \toprule
         \multirow{2}{*}{Methods} & \multirow{2}{*}{\makecell[c]{PT \\ Epochs}} & \multicolumn{2}{c}{IN1K FT} & \multicolumn{4}{c}{Fine-tuning (FT)} & \multicolumn{4}{c}{Linear Probing (LP)} \\
         \cmidrule{3-4} \cmidrule(lr){5-8} \cmidrule(lr){9-12} 
         & & Epochs & Acc & SODA & MFNet-T & SCUT-Seg & Mean & SODA & MFNet-T & SCUT-Seg & Mean \\
         \midrule
         \textcolor{gray}{ViT-Tiny/16} & & &  & & & & & & & & \\
         Scratch & - & - & - & 31.34 & 19.50 & 41.09 & 30.64 & - & - & - & - \\
         \rowcolor{gray!15} MAE$^\dag$ \citep{mae} & 800 & 200 & \underline{71.8} & 52.85 & 35.93 & 51.31 & 46.70 & 23.75 & 15.79 & 27.18 & 22.24 \\
         \rowcolor{orange!15} DeiT \citep{deit} & 300 & - & \textbf{72.2} & 63.14 & 44.60 & 61.36 & 56.37 & 42.29 & 21.78 & 31.96 & 32.01 \\
         \rowcolor{cyan!15} UNIP (MAE-L) & 100 & - & - & \underline{64.83} & \textbf{48.77} & \underline{67.22} & \underline{60.27} & \underline{44.12} & \underline{28.26} & \underline{35.09} & \underline{35.82} \\
         \rowcolor{cyan!15} UNIP (iBOT-L) & 100 & - & - & \textbf{65.54} & \underline{48.45} & \textbf{67.73} & \textbf{60.57} & \textbf{52.95} & \textbf{30.10} & \textbf{40.12} & \textbf{41.06}  \\
         \midrule
         \textcolor{gray}{ViT-Small/16} & & & & & & & & & & & \\
         Scratch & - & - & - & 41.70 & 22.49 & 46.28 & 36.82 & - & - & - & - \\
         \rowcolor{gray!15} MAE$^\dag$ \citep{mae} & 800 & 200 & 80.0 & 63.36 & 42.44 & 60.38 & 55.39 & 38.17 & 21.14 & 34.15 & 31.15 \\
         \rowcolor{gray!15} CrossMAE \citep{crossmae} & 800 & 200 & 80.5 & 63.95 & 43.99 & 63.53 & 57.16 & 39.40 & 23.87 & 34.01 & 32.43 \\
         \rowcolor{orange!15} DeiT \citep{deit} & 300 & - & 79.9 & 68.08 & 45.91 & 66.17 & 60.05 & 44.88 & 28.53 & 38.92 & 37.44 \\
         \rowcolor{orange!15} DeiT III \citep{deit3} & 800 & - & 81.4 & 69.35 & 47.73 & 67.32 & 61.47 & 54.17 & 32.01 & 43.54 & 43.24 \\
         \rowcolor{orange!15} DINO \citep{dino} & 3200$^\star$ & 200 & \underline{82.0} & 68.56 & 47.98 & 68.74 & 61.76 & 56.02 & 32.94 & 45.94 & 44.97 \\
         \rowcolor{orange!15} iBOT \citep{iBOT} & 3200$^\star$ & 200 & \textbf{82.3} & 69.33 & 47.15 & 69.80 & 62.09 & 57.10 & 33.87 & 45.82 & 45.60 \\
         \rowcolor{cyan!15} UNIP (DINO-B) & 100 & - & - & 69.35 & 49.95 & 69.70 & 63.00 & \underline{57.76} & \underline{34.15} & \underline{46.37} & \underline{46.09} \\
         \rowcolor{cyan!15} UNIP (MAE-L) & 100 & - & - & \textbf{70.99} & \underline{51.32} & \underline{70.79} & \underline{64.37} & 55.25 & 33.49 & 43.37 & 44.04 \\
         \rowcolor{cyan!15} UNIP (iBOT-L) & 100 & - & - & \underline{70.75} & \textbf{51.81} & \textbf{71.55} & \textbf{64.70} & \textbf{60.28} & \textbf{37.16} & \textbf{47.68} & \textbf{48.37} \\ 
        \midrule
        \textcolor{gray}{ViT-Base/16} & & & & & & & & & & & \\
        Scratch & - & - & - & 44.25 & 23.72 & 49.44 & 39.14 & - & - & - & - \\
        \rowcolor{gray!15} MAE \citep{mae} & 1600 & 100 & 83.6 & 68.18 & 46.78 & 67.86 & 60.94 & 43.01 & 23.42 & 37.48 & 34.64 \\
        \rowcolor{gray!15} CrossMAE \citep{crossmae} & 800 & 100 & 83.7 & 68.29 & 47.85 & 68.39 & 61.51 & 43.35 & 26.03 & 38.36 & 35.91 \\
        \rowcolor{orange!15} DeiT \citep{deit} & 300 & - & 81.8 & 69.73 & 48.59 & 69.35 & 62.56 & 57.40 & 34.82 & 46.44 & 46.22 \\
        \rowcolor{orange!15} DeiT III \citep{deit3} & 800 & 20 & \underline{83.8} & 71.09 & 49.62 & 70.19 & 63.63 & 59.01 & \underline{35.34} & 48.01 & 47.45 \\
        \rowcolor{orange!15} DINO \citep{dino} & 1600$^\star$ & 100 & 83.6 & 69.79 & 48.54 & 69.82 & 62.72 & 59.33 & 34.86 & 47.23 & 47.14 \\
        \rowcolor{orange!15} iBOT \citep{iBOT} & 1600$^\star$ & 100 & \textbf{84.0} & 71.15 & 48.98 & 71.26 & 63.80 & \underline{60.05} & 34.34 & \underline{49.12} & \underline{47.84} \\
        \rowcolor{cyan!15} UNIP (MAE-L) & 100 & - & - & \underline{71.47} & \textbf{52.55} & \underline{71.82} & \textbf{65.28} & 58.82 & 34.75 & 48.74 & 47.43 \\
        \rowcolor{cyan!15} UNIP (iBOT-L) & 100 & - & - & \textbf{71.75} & \underline{51.46} & \textbf{72.00} & \underline{65.07} & \textbf{63.14} & \textbf{39.08} & \textbf{52.53} & \textbf{51.58} \\
        \midrule
        \textcolor{gray}{ViT-Large/16} & & & & & & & & & & & \\
        Scratch & - & - & - & 44.70 & 23.68 & 49.55 & 39.31 & - & - & - & - \\
        \rowcolor{gray!15} MAE \citep{mae} & 1600 & 50 & \textbf{85.9} & 71.04 & \underline{51.17} & 70.83 & 64.35 & 52.20 & 31.21 & 43.71 & 42.37 \\
        \rowcolor{gray!15} CrossMAE \citep{crossmae} & 800 & 50 & 85.4 & 70.48 & 50.97 & 70.24 & 63.90 & 53.29 & 33.09 & 45.01 & 43.80 \\
        \rowcolor{orange!15} DeiT3 \citep{deit3} & 800 & 20 & \underline{84.9} & \underline{71.67} & 50.78 & \textbf{71.54} & \underline{64.66} & \underline{59.42} & \textbf{37.57} & \textbf{50.27} & \underline{49.09} \\
        \rowcolor{orange!15} iBOT \citep{iBOT} & 1000$^\star$ & 50 & 84.8 & \textbf{71.75} & \textbf{51.66} & \underline{71.49} & \textbf{64.97} & \textbf{61.73} & \underline{36.68} & \underline{50.12} & \textbf{49.51} \\
        \bottomrule
    \end{tabular}}}
    \vspace{-2mm}
\end{table}

\section{Future Directions}
Based on our experimental findings, we propose two potential research directions that integrate role-playing and reasoning techniques for future exploration.

\noindent \paragraph{Role-aware Chain-of-Thought (CoT) for Improving Role-playing LLMs.}
One promising direction for enhancing the role-playing ability of LLMs is Role-aware CoT reasoning. 
While standard CoT enables step-by-step logical inference, it often disregards the persona-specific constraints that are essential for consistent role-playing. 
A role-aware CoT approach would integrate persona attributes, narrative constraints, and character-specific perspectives into the reasoning process, ensuring that logical deductions align with the character's predefined traits. 
For instance, a historical figure simulated in an LLM should reason within the knowledge and biases of their era rather than applying contemporary logic. 
This requires incorporating dynamic memory structures that retain persona information throughout multi-turn interactions, mitigating the risk of breaking character or adopting inconsistent reasoning patterns.

\noindent\paragraph{Reinforcement Learning for Role-playing LLMs.}

DeepSeek-R1 has demonstrated that by defining precise rule-based rewards, reinforcement learning alone can induce emergent reasoning and cognitive capabilities. This suggests an important research direction: investigating whether carefully designed role-playing task rewards can enable models to autonomously develop intrinsic, role-specific reasoning and thinking abilities, thereby enhancing their role-playing performance.
A critical consideration in reward design is ensuring that it simultaneously accounts for the accuracy of the role-playing task while effectively guiding the model to develop role-specific reasoning abilities. The reward function should prevent the model from relying on shortcuts to generate final responses, instead encouraging it to engage in authentic character-driven reasoning and decision-making.

\section{Conclusion}
In this work, we propose a simple yet effective approach, called SMILE, for graph few-shot learning with fewer tasks. Specifically, we introduce a novel dual-level mixup strategy, including within-task and across-task mixup, for enriching the diversity of nodes within each task and the diversity of tasks. Also, we incorporate the degree-based prior information to learn expressive node embeddings. Theoretically, we prove that SMILE effectively enhances the model's generalization performance. Empirically, we conduct extensive experiments on multiple benchmarks and the results suggest that SMILE significantly outperforms other baselines, including both in-domain and cross-domain few-shot settings.


\bibliography{reference}
% \bibliography{new_reference}
\bibliographystyle{IEEEtran}

\subsection{Lloyd-Max Algorithm}
\label{subsec:Lloyd-Max}
For a given quantization bitwidth $B$ and an operand $\bm{X}$, the Lloyd-Max algorithm finds $2^B$ quantization levels $\{\hat{x}_i\}_{i=1}^{2^B}$ such that quantizing $\bm{X}$ by rounding each scalar in $\bm{X}$ to the nearest quantization level minimizes the quantization MSE. 

The algorithm starts with an initial guess of quantization levels and then iteratively computes quantization thresholds $\{\tau_i\}_{i=1}^{2^B-1}$ and updates quantization levels $\{\hat{x}_i\}_{i=1}^{2^B}$. Specifically, at iteration $n$, thresholds are set to the midpoints of the previous iteration's levels:
\begin{align*}
    \tau_i^{(n)}=\frac{\hat{x}_i^{(n-1)}+\hat{x}_{i+1}^{(n-1)}}2 \text{ for } i=1\ldots 2^B-1
\end{align*}
Subsequently, the quantization levels are re-computed as conditional means of the data regions defined by the new thresholds:
\begin{align*}
    \hat{x}_i^{(n)}=\mathbb{E}\left[ \bm{X} \big| \bm{X}\in [\tau_{i-1}^{(n)},\tau_i^{(n)}] \right] \text{ for } i=1\ldots 2^B
\end{align*}
where to satisfy boundary conditions we have $\tau_0=-\infty$ and $\tau_{2^B}=\infty$. The algorithm iterates the above steps until convergence.

Figure \ref{fig:lm_quant} compares the quantization levels of a $7$-bit floating point (E3M3) quantizer (left) to a $7$-bit Lloyd-Max quantizer (right) when quantizing a layer of weights from the GPT3-126M model at a per-tensor granularity. As shown, the Lloyd-Max quantizer achieves substantially lower quantization MSE. Further, Table \ref{tab:FP7_vs_LM7} shows the superior perplexity achieved by Lloyd-Max quantizers for bitwidths of $7$, $6$ and $5$. The difference between the quantizers is clear at 5 bits, where per-tensor FP quantization incurs a drastic and unacceptable increase in perplexity, while Lloyd-Max quantization incurs a much smaller increase. Nevertheless, we note that even the optimal Lloyd-Max quantizer incurs a notable ($\sim 1.5$) increase in perplexity due to the coarse granularity of quantization. 

\begin{figure}[h]
  \centering
  \includegraphics[width=0.7\linewidth]{sections/figures/LM7_FP7.pdf}
  \caption{\small Quantization levels and the corresponding quantization MSE of Floating Point (left) vs Lloyd-Max (right) Quantizers for a layer of weights in the GPT3-126M model.}
  \label{fig:lm_quant}
\end{figure}

\begin{table}[h]\scriptsize
\begin{center}
\caption{\label{tab:FP7_vs_LM7} \small Comparing perplexity (lower is better) achieved by floating point quantizers and Lloyd-Max quantizers on a GPT3-126M model for the Wikitext-103 dataset.}
\begin{tabular}{c|cc|c}
\hline
 \multirow{2}{*}{\textbf{Bitwidth}} & \multicolumn{2}{|c|}{\textbf{Floating-Point Quantizer}} & \textbf{Lloyd-Max Quantizer} \\
 & Best Format & Wikitext-103 Perplexity & Wikitext-103 Perplexity \\
\hline
7 & E3M3 & 18.32 & 18.27 \\
6 & E3M2 & 19.07 & 18.51 \\
5 & E4M0 & 43.89 & 19.71 \\
\hline
\end{tabular}
\end{center}
\end{table}

\subsection{Proof of Local Optimality of LO-BCQ}
\label{subsec:lobcq_opt_proof}
For a given block $\bm{b}_j$, the quantization MSE during LO-BCQ can be empirically evaluated as $\frac{1}{L_b}\lVert \bm{b}_j- \bm{\hat{b}}_j\rVert^2_2$ where $\bm{\hat{b}}_j$ is computed from equation (\ref{eq:clustered_quantization_definition}) as $C_{f(\bm{b}_j)}(\bm{b}_j)$. Further, for a given block cluster $\mathcal{B}_i$, we compute the quantization MSE as $\frac{1}{|\mathcal{B}_{i}|}\sum_{\bm{b} \in \mathcal{B}_{i}} \frac{1}{L_b}\lVert \bm{b}- C_i^{(n)}(\bm{b})\rVert^2_2$. Therefore, at the end of iteration $n$, we evaluate the overall quantization MSE $J^{(n)}$ for a given operand $\bm{X}$ composed of $N_c$ block clusters as:
\begin{align*}
    \label{eq:mse_iter_n}
    J^{(n)} = \frac{1}{N_c} \sum_{i=1}^{N_c} \frac{1}{|\mathcal{B}_{i}^{(n)}|}\sum_{\bm{v} \in \mathcal{B}_{i}^{(n)}} \frac{1}{L_b}\lVert \bm{b}- B_i^{(n)}(\bm{b})\rVert^2_2
\end{align*}

At the end of iteration $n$, the codebooks are updated from $\mathcal{C}^{(n-1)}$ to $\mathcal{C}^{(n)}$. However, the mapping of a given vector $\bm{b}_j$ to quantizers $\mathcal{C}^{(n)}$ remains as  $f^{(n)}(\bm{b}_j)$. At the next iteration, during the vector clustering step, $f^{(n+1)}(\bm{b}_j)$ finds new mapping of $\bm{b}_j$ to updated codebooks $\mathcal{C}^{(n)}$ such that the quantization MSE over the candidate codebooks is minimized. Therefore, we obtain the following result for $\bm{b}_j$:
\begin{align*}
\frac{1}{L_b}\lVert \bm{b}_j - C_{f^{(n+1)}(\bm{b}_j)}^{(n)}(\bm{b}_j)\rVert^2_2 \le \frac{1}{L_b}\lVert \bm{b}_j - C_{f^{(n)}(\bm{b}_j)}^{(n)}(\bm{b}_j)\rVert^2_2
\end{align*}

That is, quantizing $\bm{b}_j$ at the end of the block clustering step of iteration $n+1$ results in lower quantization MSE compared to quantizing at the end of iteration $n$. Since this is true for all $\bm{b} \in \bm{X}$, we assert the following:
\begin{equation}
\begin{split}
\label{eq:mse_ineq_1}
    \tilde{J}^{(n+1)} &= \frac{1}{N_c} \sum_{i=1}^{N_c} \frac{1}{|\mathcal{B}_{i}^{(n+1)}|}\sum_{\bm{b} \in \mathcal{B}_{i}^{(n+1)}} \frac{1}{L_b}\lVert \bm{b} - C_i^{(n)}(b)\rVert^2_2 \le J^{(n)}
\end{split}
\end{equation}
where $\tilde{J}^{(n+1)}$ is the the quantization MSE after the vector clustering step at iteration $n+1$.

Next, during the codebook update step (\ref{eq:quantizers_update}) at iteration $n+1$, the per-cluster codebooks $\mathcal{C}^{(n)}$ are updated to $\mathcal{C}^{(n+1)}$ by invoking the Lloyd-Max algorithm \citep{Lloyd}. We know that for any given value distribution, the Lloyd-Max algorithm minimizes the quantization MSE. Therefore, for a given vector cluster $\mathcal{B}_i$ we obtain the following result:

\begin{equation}
    \frac{1}{|\mathcal{B}_{i}^{(n+1)}|}\sum_{\bm{b} \in \mathcal{B}_{i}^{(n+1)}} \frac{1}{L_b}\lVert \bm{b}- C_i^{(n+1)}(\bm{b})\rVert^2_2 \le \frac{1}{|\mathcal{B}_{i}^{(n+1)}|}\sum_{\bm{b} \in \mathcal{B}_{i}^{(n+1)}} \frac{1}{L_b}\lVert \bm{b}- C_i^{(n)}(\bm{b})\rVert^2_2
\end{equation}

The above equation states that quantizing the given block cluster $\mathcal{B}_i$ after updating the associated codebook from $C_i^{(n)}$ to $C_i^{(n+1)}$ results in lower quantization MSE. Since this is true for all the block clusters, we derive the following result: 
\begin{equation}
\begin{split}
\label{eq:mse_ineq_2}
     J^{(n+1)} &= \frac{1}{N_c} \sum_{i=1}^{N_c} \frac{1}{|\mathcal{B}_{i}^{(n+1)}|}\sum_{\bm{b} \in \mathcal{B}_{i}^{(n+1)}} \frac{1}{L_b}\lVert \bm{b}- C_i^{(n+1)}(\bm{b})\rVert^2_2  \le \tilde{J}^{(n+1)}   
\end{split}
\end{equation}

Following (\ref{eq:mse_ineq_1}) and (\ref{eq:mse_ineq_2}), we find that the quantization MSE is non-increasing for each iteration, that is, $J^{(1)} \ge J^{(2)} \ge J^{(3)} \ge \ldots \ge J^{(M)}$ where $M$ is the maximum number of iterations. 
%Therefore, we can say that if the algorithm converges, then it must be that it has converged to a local minimum. 
\hfill $\blacksquare$


\begin{figure}
    \begin{center}
    \includegraphics[width=0.5\textwidth]{sections//figures/mse_vs_iter.pdf}
    \end{center}
    \caption{\small NMSE vs iterations during LO-BCQ compared to other block quantization proposals}
    \label{fig:nmse_vs_iter}
\end{figure}

Figure \ref{fig:nmse_vs_iter} shows the empirical convergence of LO-BCQ across several block lengths and number of codebooks. Also, the MSE achieved by LO-BCQ is compared to baselines such as MXFP and VSQ. As shown, LO-BCQ converges to a lower MSE than the baselines. Further, we achieve better convergence for larger number of codebooks ($N_c$) and for a smaller block length ($L_b$), both of which increase the bitwidth of BCQ (see Eq \ref{eq:bitwidth_bcq}).


\subsection{Additional Accuracy Results}
%Table \ref{tab:lobcq_config} lists the various LOBCQ configurations and their corresponding bitwidths.
\begin{table}
\setlength{\tabcolsep}{4.75pt}
\begin{center}
\caption{\label{tab:lobcq_config} Various LO-BCQ configurations and their bitwidths.}
\begin{tabular}{|c||c|c|c|c||c|c||c|} 
\hline
 & \multicolumn{4}{|c||}{$L_b=8$} & \multicolumn{2}{|c||}{$L_b=4$} & $L_b=2$ \\
 \hline
 \backslashbox{$L_A$\kern-1em}{\kern-1em$N_c$} & 2 & 4 & 8 & 16 & 2 & 4 & 2 \\
 \hline
 64 & 4.25 & 4.375 & 4.5 & 4.625 & 4.375 & 4.625 & 4.625\\
 \hline
 32 & 4.375 & 4.5 & 4.625& 4.75 & 4.5 & 4.75 & 4.75 \\
 \hline
 16 & 4.625 & 4.75& 4.875 & 5 & 4.75 & 5 & 5 \\
 \hline
\end{tabular}
\end{center}
\end{table}

%\subsection{Perplexity achieved by various LO-BCQ configurations on Wikitext-103 dataset}

\begin{table} \centering
\begin{tabular}{|c||c|c|c|c||c|c||c|} 
\hline
 $L_b \rightarrow$& \multicolumn{4}{c||}{8} & \multicolumn{2}{c||}{4} & 2\\
 \hline
 \backslashbox{$L_A$\kern-1em}{\kern-1em$N_c$} & 2 & 4 & 8 & 16 & 2 & 4 & 2  \\
 %$N_c \rightarrow$ & 2 & 4 & 8 & 16 & 2 & 4 & 2 \\
 \hline
 \hline
 \multicolumn{8}{c}{GPT3-1.3B (FP32 PPL = 9.98)} \\ 
 \hline
 \hline
 64 & 10.40 & 10.23 & 10.17 & 10.15 &  10.28 & 10.18 & 10.19 \\
 \hline
 32 & 10.25 & 10.20 & 10.15 & 10.12 &  10.23 & 10.17 & 10.17 \\
 \hline
 16 & 10.22 & 10.16 & 10.10 & 10.09 &  10.21 & 10.14 & 10.16 \\
 \hline
  \hline
 \multicolumn{8}{c}{GPT3-8B (FP32 PPL = 7.38)} \\ 
 \hline
 \hline
 64 & 7.61 & 7.52 & 7.48 &  7.47 &  7.55 &  7.49 & 7.50 \\
 \hline
 32 & 7.52 & 7.50 & 7.46 &  7.45 &  7.52 &  7.48 & 7.48  \\
 \hline
 16 & 7.51 & 7.48 & 7.44 &  7.44 &  7.51 &  7.49 & 7.47  \\
 \hline
\end{tabular}
\caption{\label{tab:ppl_gpt3_abalation} Wikitext-103 perplexity across GPT3-1.3B and 8B models.}
\end{table}

\begin{table} \centering
\begin{tabular}{|c||c|c|c|c||} 
\hline
 $L_b \rightarrow$& \multicolumn{4}{c||}{8}\\
 \hline
 \backslashbox{$L_A$\kern-1em}{\kern-1em$N_c$} & 2 & 4 & 8 & 16 \\
 %$N_c \rightarrow$ & 2 & 4 & 8 & 16 & 2 & 4 & 2 \\
 \hline
 \hline
 \multicolumn{5}{|c|}{Llama2-7B (FP32 PPL = 5.06)} \\ 
 \hline
 \hline
 64 & 5.31 & 5.26 & 5.19 & 5.18  \\
 \hline
 32 & 5.23 & 5.25 & 5.18 & 5.15  \\
 \hline
 16 & 5.23 & 5.19 & 5.16 & 5.14  \\
 \hline
 \multicolumn{5}{|c|}{Nemotron4-15B (FP32 PPL = 5.87)} \\ 
 \hline
 \hline
 64  & 6.3 & 6.20 & 6.13 & 6.08  \\
 \hline
 32  & 6.24 & 6.12 & 6.07 & 6.03  \\
 \hline
 16  & 6.12 & 6.14 & 6.04 & 6.02  \\
 \hline
 \multicolumn{5}{|c|}{Nemotron4-340B (FP32 PPL = 3.48)} \\ 
 \hline
 \hline
 64 & 3.67 & 3.62 & 3.60 & 3.59 \\
 \hline
 32 & 3.63 & 3.61 & 3.59 & 3.56 \\
 \hline
 16 & 3.61 & 3.58 & 3.57 & 3.55 \\
 \hline
\end{tabular}
\caption{\label{tab:ppl_llama7B_nemo15B} Wikitext-103 perplexity compared to FP32 baseline in Llama2-7B and Nemotron4-15B, 340B models}
\end{table}

%\subsection{Perplexity achieved by various LO-BCQ configurations on MMLU dataset}


\begin{table} \centering
\begin{tabular}{|c||c|c|c|c||c|c|c|c|} 
\hline
 $L_b \rightarrow$& \multicolumn{4}{c||}{8} & \multicolumn{4}{c||}{8}\\
 \hline
 \backslashbox{$L_A$\kern-1em}{\kern-1em$N_c$} & 2 & 4 & 8 & 16 & 2 & 4 & 8 & 16  \\
 %$N_c \rightarrow$ & 2 & 4 & 8 & 16 & 2 & 4 & 2 \\
 \hline
 \hline
 \multicolumn{5}{|c|}{Llama2-7B (FP32 Accuracy = 45.8\%)} & \multicolumn{4}{|c|}{Llama2-70B (FP32 Accuracy = 69.12\%)} \\ 
 \hline
 \hline
 64 & 43.9 & 43.4 & 43.9 & 44.9 & 68.07 & 68.27 & 68.17 & 68.75 \\
 \hline
 32 & 44.5 & 43.8 & 44.9 & 44.5 & 68.37 & 68.51 & 68.35 & 68.27  \\
 \hline
 16 & 43.9 & 42.7 & 44.9 & 45 & 68.12 & 68.77 & 68.31 & 68.59  \\
 \hline
 \hline
 \multicolumn{5}{|c|}{GPT3-22B (FP32 Accuracy = 38.75\%)} & \multicolumn{4}{|c|}{Nemotron4-15B (FP32 Accuracy = 64.3\%)} \\ 
 \hline
 \hline
 64 & 36.71 & 38.85 & 38.13 & 38.92 & 63.17 & 62.36 & 63.72 & 64.09 \\
 \hline
 32 & 37.95 & 38.69 & 39.45 & 38.34 & 64.05 & 62.30 & 63.8 & 64.33  \\
 \hline
 16 & 38.88 & 38.80 & 38.31 & 38.92 & 63.22 & 63.51 & 63.93 & 64.43  \\
 \hline
\end{tabular}
\caption{\label{tab:mmlu_abalation} Accuracy on MMLU dataset across GPT3-22B, Llama2-7B, 70B and Nemotron4-15B models.}
\end{table}


%\subsection{Perplexity achieved by various LO-BCQ configurations on LM evaluation harness}

\begin{table} \centering
\begin{tabular}{|c||c|c|c|c||c|c|c|c|} 
\hline
 $L_b \rightarrow$& \multicolumn{4}{c||}{8} & \multicolumn{4}{c||}{8}\\
 \hline
 \backslashbox{$L_A$\kern-1em}{\kern-1em$N_c$} & 2 & 4 & 8 & 16 & 2 & 4 & 8 & 16  \\
 %$N_c \rightarrow$ & 2 & 4 & 8 & 16 & 2 & 4 & 2 \\
 \hline
 \hline
 \multicolumn{5}{|c|}{Race (FP32 Accuracy = 37.51\%)} & \multicolumn{4}{|c|}{Boolq (FP32 Accuracy = 64.62\%)} \\ 
 \hline
 \hline
 64 & 36.94 & 37.13 & 36.27 & 37.13 & 63.73 & 62.26 & 63.49 & 63.36 \\
 \hline
 32 & 37.03 & 36.36 & 36.08 & 37.03 & 62.54 & 63.51 & 63.49 & 63.55  \\
 \hline
 16 & 37.03 & 37.03 & 36.46 & 37.03 & 61.1 & 63.79 & 63.58 & 63.33  \\
 \hline
 \hline
 \multicolumn{5}{|c|}{Winogrande (FP32 Accuracy = 58.01\%)} & \multicolumn{4}{|c|}{Piqa (FP32 Accuracy = 74.21\%)} \\ 
 \hline
 \hline
 64 & 58.17 & 57.22 & 57.85 & 58.33 & 73.01 & 73.07 & 73.07 & 72.80 \\
 \hline
 32 & 59.12 & 58.09 & 57.85 & 58.41 & 73.01 & 73.94 & 72.74 & 73.18  \\
 \hline
 16 & 57.93 & 58.88 & 57.93 & 58.56 & 73.94 & 72.80 & 73.01 & 73.94  \\
 \hline
\end{tabular}
\caption{\label{tab:mmlu_abalation} Accuracy on LM evaluation harness tasks on GPT3-1.3B model.}
\end{table}

\begin{table} \centering
\begin{tabular}{|c||c|c|c|c||c|c|c|c|} 
\hline
 $L_b \rightarrow$& \multicolumn{4}{c||}{8} & \multicolumn{4}{c||}{8}\\
 \hline
 \backslashbox{$L_A$\kern-1em}{\kern-1em$N_c$} & 2 & 4 & 8 & 16 & 2 & 4 & 8 & 16  \\
 %$N_c \rightarrow$ & 2 & 4 & 8 & 16 & 2 & 4 & 2 \\
 \hline
 \hline
 \multicolumn{5}{|c|}{Race (FP32 Accuracy = 41.34\%)} & \multicolumn{4}{|c|}{Boolq (FP32 Accuracy = 68.32\%)} \\ 
 \hline
 \hline
 64 & 40.48 & 40.10 & 39.43 & 39.90 & 69.20 & 68.41 & 69.45 & 68.56 \\
 \hline
 32 & 39.52 & 39.52 & 40.77 & 39.62 & 68.32 & 67.43 & 68.17 & 69.30  \\
 \hline
 16 & 39.81 & 39.71 & 39.90 & 40.38 & 68.10 & 66.33 & 69.51 & 69.42  \\
 \hline
 \hline
 \multicolumn{5}{|c|}{Winogrande (FP32 Accuracy = 67.88\%)} & \multicolumn{4}{|c|}{Piqa (FP32 Accuracy = 78.78\%)} \\ 
 \hline
 \hline
 64 & 66.85 & 66.61 & 67.72 & 67.88 & 77.31 & 77.42 & 77.75 & 77.64 \\
 \hline
 32 & 67.25 & 67.72 & 67.72 & 67.00 & 77.31 & 77.04 & 77.80 & 77.37  \\
 \hline
 16 & 68.11 & 68.90 & 67.88 & 67.48 & 77.37 & 78.13 & 78.13 & 77.69  \\
 \hline
\end{tabular}
\caption{\label{tab:mmlu_abalation} Accuracy on LM evaluation harness tasks on GPT3-8B model.}
\end{table}

\begin{table} \centering
\begin{tabular}{|c||c|c|c|c||c|c|c|c|} 
\hline
 $L_b \rightarrow$& \multicolumn{4}{c||}{8} & \multicolumn{4}{c||}{8}\\
 \hline
 \backslashbox{$L_A$\kern-1em}{\kern-1em$N_c$} & 2 & 4 & 8 & 16 & 2 & 4 & 8 & 16  \\
 %$N_c \rightarrow$ & 2 & 4 & 8 & 16 & 2 & 4 & 2 \\
 \hline
 \hline
 \multicolumn{5}{|c|}{Race (FP32 Accuracy = 40.67\%)} & \multicolumn{4}{|c|}{Boolq (FP32 Accuracy = 76.54\%)} \\ 
 \hline
 \hline
 64 & 40.48 & 40.10 & 39.43 & 39.90 & 75.41 & 75.11 & 77.09 & 75.66 \\
 \hline
 32 & 39.52 & 39.52 & 40.77 & 39.62 & 76.02 & 76.02 & 75.96 & 75.35  \\
 \hline
 16 & 39.81 & 39.71 & 39.90 & 40.38 & 75.05 & 73.82 & 75.72 & 76.09  \\
 \hline
 \hline
 \multicolumn{5}{|c|}{Winogrande (FP32 Accuracy = 70.64\%)} & \multicolumn{4}{|c|}{Piqa (FP32 Accuracy = 79.16\%)} \\ 
 \hline
 \hline
 64 & 69.14 & 70.17 & 70.17 & 70.56 & 78.24 & 79.00 & 78.62 & 78.73 \\
 \hline
 32 & 70.96 & 69.69 & 71.27 & 69.30 & 78.56 & 79.49 & 79.16 & 78.89  \\
 \hline
 16 & 71.03 & 69.53 & 69.69 & 70.40 & 78.13 & 79.16 & 79.00 & 79.00  \\
 \hline
\end{tabular}
\caption{\label{tab:mmlu_abalation} Accuracy on LM evaluation harness tasks on GPT3-22B model.}
\end{table}

\begin{table} \centering
\begin{tabular}{|c||c|c|c|c||c|c|c|c|} 
\hline
 $L_b \rightarrow$& \multicolumn{4}{c||}{8} & \multicolumn{4}{c||}{8}\\
 \hline
 \backslashbox{$L_A$\kern-1em}{\kern-1em$N_c$} & 2 & 4 & 8 & 16 & 2 & 4 & 8 & 16  \\
 %$N_c \rightarrow$ & 2 & 4 & 8 & 16 & 2 & 4 & 2 \\
 \hline
 \hline
 \multicolumn{5}{|c|}{Race (FP32 Accuracy = 44.4\%)} & \multicolumn{4}{|c|}{Boolq (FP32 Accuracy = 79.29\%)} \\ 
 \hline
 \hline
 64 & 42.49 & 42.51 & 42.58 & 43.45 & 77.58 & 77.37 & 77.43 & 78.1 \\
 \hline
 32 & 43.35 & 42.49 & 43.64 & 43.73 & 77.86 & 75.32 & 77.28 & 77.86  \\
 \hline
 16 & 44.21 & 44.21 & 43.64 & 42.97 & 78.65 & 77 & 76.94 & 77.98  \\
 \hline
 \hline
 \multicolumn{5}{|c|}{Winogrande (FP32 Accuracy = 69.38\%)} & \multicolumn{4}{|c|}{Piqa (FP32 Accuracy = 78.07\%)} \\ 
 \hline
 \hline
 64 & 68.9 & 68.43 & 69.77 & 68.19 & 77.09 & 76.82 & 77.09 & 77.86 \\
 \hline
 32 & 69.38 & 68.51 & 68.82 & 68.90 & 78.07 & 76.71 & 78.07 & 77.86  \\
 \hline
 16 & 69.53 & 67.09 & 69.38 & 68.90 & 77.37 & 77.8 & 77.91 & 77.69  \\
 \hline
\end{tabular}
\caption{\label{tab:mmlu_abalation} Accuracy on LM evaluation harness tasks on Llama2-7B model.}
\end{table}

\begin{table} \centering
\begin{tabular}{|c||c|c|c|c||c|c|c|c|} 
\hline
 $L_b \rightarrow$& \multicolumn{4}{c||}{8} & \multicolumn{4}{c||}{8}\\
 \hline
 \backslashbox{$L_A$\kern-1em}{\kern-1em$N_c$} & 2 & 4 & 8 & 16 & 2 & 4 & 8 & 16  \\
 %$N_c \rightarrow$ & 2 & 4 & 8 & 16 & 2 & 4 & 2 \\
 \hline
 \hline
 \multicolumn{5}{|c|}{Race (FP32 Accuracy = 48.8\%)} & \multicolumn{4}{|c|}{Boolq (FP32 Accuracy = 85.23\%)} \\ 
 \hline
 \hline
 64 & 49.00 & 49.00 & 49.28 & 48.71 & 82.82 & 84.28 & 84.03 & 84.25 \\
 \hline
 32 & 49.57 & 48.52 & 48.33 & 49.28 & 83.85 & 84.46 & 84.31 & 84.93  \\
 \hline
 16 & 49.85 & 49.09 & 49.28 & 48.99 & 85.11 & 84.46 & 84.61 & 83.94  \\
 \hline
 \hline
 \multicolumn{5}{|c|}{Winogrande (FP32 Accuracy = 79.95\%)} & \multicolumn{4}{|c|}{Piqa (FP32 Accuracy = 81.56\%)} \\ 
 \hline
 \hline
 64 & 78.77 & 78.45 & 78.37 & 79.16 & 81.45 & 80.69 & 81.45 & 81.5 \\
 \hline
 32 & 78.45 & 79.01 & 78.69 & 80.66 & 81.56 & 80.58 & 81.18 & 81.34  \\
 \hline
 16 & 79.95 & 79.56 & 79.79 & 79.72 & 81.28 & 81.66 & 81.28 & 80.96  \\
 \hline
\end{tabular}
\caption{\label{tab:mmlu_abalation} Accuracy on LM evaluation harness tasks on Llama2-70B model.}
\end{table}

%\section{MSE Studies}
%\textcolor{red}{TODO}


\subsection{Number Formats and Quantization Method}
\label{subsec:numFormats_quantMethod}
\subsubsection{Integer Format}
An $n$-bit signed integer (INT) is typically represented with a 2s-complement format \citep{yao2022zeroquant,xiao2023smoothquant,dai2021vsq}, where the most significant bit denotes the sign.

\subsubsection{Floating Point Format}
An $n$-bit signed floating point (FP) number $x$ comprises of a 1-bit sign ($x_{\mathrm{sign}}$), $B_m$-bit mantissa ($x_{\mathrm{mant}}$) and $B_e$-bit exponent ($x_{\mathrm{exp}}$) such that $B_m+B_e=n-1$. The associated constant exponent bias ($E_{\mathrm{bias}}$) is computed as $(2^{{B_e}-1}-1)$. We denote this format as $E_{B_e}M_{B_m}$.  

\subsubsection{Quantization Scheme}
\label{subsec:quant_method}
A quantization scheme dictates how a given unquantized tensor is converted to its quantized representation. We consider FP formats for the purpose of illustration. Given an unquantized tensor $\bm{X}$ and an FP format $E_{B_e}M_{B_m}$, we first, we compute the quantization scale factor $s_X$ that maps the maximum absolute value of $\bm{X}$ to the maximum quantization level of the $E_{B_e}M_{B_m}$ format as follows:
\begin{align}
\label{eq:sf}
    s_X = \frac{\mathrm{max}(|\bm{X}|)}{\mathrm{max}(E_{B_e}M_{B_m})}
\end{align}
In the above equation, $|\cdot|$ denotes the absolute value function.

Next, we scale $\bm{X}$ by $s_X$ and quantize it to $\hat{\bm{X}}$ by rounding it to the nearest quantization level of $E_{B_e}M_{B_m}$ as:

\begin{align}
\label{eq:tensor_quant}
    \hat{\bm{X}} = \text{round-to-nearest}\left(\frac{\bm{X}}{s_X}, E_{B_e}M_{B_m}\right)
\end{align}

We perform dynamic max-scaled quantization \citep{wu2020integer}, where the scale factor $s$ for activations is dynamically computed during runtime.

\subsection{Vector Scaled Quantization}
\begin{wrapfigure}{r}{0.35\linewidth}
  \centering
  \includegraphics[width=\linewidth]{sections/figures/vsquant.jpg}
  \caption{\small Vectorwise decomposition for per-vector scaled quantization (VSQ \citep{dai2021vsq}).}
  \label{fig:vsquant}
\end{wrapfigure}
During VSQ \citep{dai2021vsq}, the operand tensors are decomposed into 1D vectors in a hardware friendly manner as shown in Figure \ref{fig:vsquant}. Since the decomposed tensors are used as operands in matrix multiplications during inference, it is beneficial to perform this decomposition along the reduction dimension of the multiplication. The vectorwise quantization is performed similar to tensorwise quantization described in Equations \ref{eq:sf} and \ref{eq:tensor_quant}, where a scale factor $s_v$ is required for each vector $\bm{v}$ that maps the maximum absolute value of that vector to the maximum quantization level. While smaller vector lengths can lead to larger accuracy gains, the associated memory and computational overheads due to the per-vector scale factors increases. To alleviate these overheads, VSQ \citep{dai2021vsq} proposed a second level quantization of the per-vector scale factors to unsigned integers, while MX \citep{rouhani2023shared} quantizes them to integer powers of 2 (denoted as $2^{INT}$).

\subsubsection{MX Format}
The MX format proposed in \citep{rouhani2023microscaling} introduces the concept of sub-block shifting. For every two scalar elements of $b$-bits each, there is a shared exponent bit. The value of this exponent bit is determined through an empirical analysis that targets minimizing quantization MSE. We note that the FP format $E_{1}M_{b}$ is strictly better than MX from an accuracy perspective since it allocates a dedicated exponent bit to each scalar as opposed to sharing it across two scalars. Therefore, we conservatively bound the accuracy of a $b+2$-bit signed MX format with that of a $E_{1}M_{b}$ format in our comparisons. For instance, we use E1M2 format as a proxy for MX4.

\begin{figure}
    \centering
    \includegraphics[width=1\linewidth]{sections//figures/BlockFormats.pdf}
    \caption{\small Comparing LO-BCQ to MX format.}
    \label{fig:block_formats}
\end{figure}

Figure \ref{fig:block_formats} compares our $4$-bit LO-BCQ block format to MX \citep{rouhani2023microscaling}. As shown, both LO-BCQ and MX decompose a given operand tensor into block arrays and each block array into blocks. Similar to MX, we find that per-block quantization ($L_b < L_A$) leads to better accuracy due to increased flexibility. While MX achieves this through per-block $1$-bit micro-scales, we associate a dedicated codebook to each block through a per-block codebook selector. Further, MX quantizes the per-block array scale-factor to E8M0 format without per-tensor scaling. In contrast during LO-BCQ, we find that per-tensor scaling combined with quantization of per-block array scale-factor to E4M3 format results in superior inference accuracy across models. 



\vfill

\end{document}


